\documentclass[reprint]{revtex4-2}
\usepackage{amsmath,amssymb,graphicx,hyperref,booktabs,enumitem}
\usepackage[colorlinks=true,citecolor=blue,linkcolor=blue]{hyperref}
\begin{document}

\title{Supplemental Material:\\
Information-Theoretic Origin of the Gravitational Potential Equation}

\author{Huang Zhongchang}
\affiliation{Independent Researcher, Nanjing, China}

\date{\today}

\maketitle
\tableofcontents

%%%%%%%%%%%%%%%%%%%%%%%%%%%%%%%%%%%%%%%%%%%%%%%%%%%%%%%%%%%%%%%%%%%%%
\section{Preliminary Notation and Conventions}
%%%%%%%%%%%%%%%%%%%%%%%%%%%%%%%%%%%%%%%%%%%%%%%%%%%%%%%%%%%%%%%%%%%%%

The main paper defines the $q$-field framework on a set of $N$ binary cells $\{s_i\in\{0,1\}\}_{i=1}^N$, with dynamics $f:X\to X$ where $X=\{0,1\}^N$. The vacant capacity fraction is $q=N^{-1}\sum_i(1-s_i)$. In the continuum limit $N\to\infty$ with cell spacing $a\to 0$, $q$ becomes a scalar field $q(\mathbf{x},t)$ on $\mathbb{R}^3$.

The flux takes the self-information (surprisal) form:
\begin{equation}
J = -\kappa\nabla(\ln q),
\end{equation}
with effective diffusivity $D(q)=\kappa/q$. The continuum $q$-field equation is:
\begin{equation}
\boxed{\partial_t q = \kappa\,\nabla^2(\ln q)}.
\end{equation}

We denote the information potential $\psi\equiv -\ln q$, so that $\partial_t q = -\kappa\nabla^2\psi$ and in the static limit $\nabla^2\psi=0$. This Supplemental Material provides complete derivations for all statements made in the main text. Every equation is numbered for traceability, and all approximations are explicitly flagged.

%%%%%%%%%%%%%%%%%%%%%%%%%%%%%%%%%%%%%%%%%%%%%%%%%%%%%%%%%%%%%%%%%%%%%
\section{Detailed Derivations}
%%%%%%%%%%%%%%%%%%%%%%%%%%%%%%%%%%%%%%%%%%%%%%%%%%%%%%%%%%%%%%%%%%%%%

% ============================================================
\subsection{A1. Proof that A1+A2 $\Rightarrow$ $f$ is a Permutation $\Rightarrow$ Cyclic Decomposition}
% ============================================================

\subsubsection{A1.1 Injectivity implies bijectivity on a finite set}

\textbf{A1 (Distinguishability):} For any two states $s\neq s'$, their images satisfy $f(s)\neq f(s')$. That is, $f$ is injective.

\textbf{A2 (Bounded capacity):} $X=\{0,1\}^N$ is finite, with $|X|=2^N$.

\begin{theorem}[Bijectivity]
\label{thm:bijectivity}
An injective map $f:X\to X$ on a finite set $X$ is bijective.
\end{theorem}

\begin{proof}
Since $f$ is injective, $|f(X)| = |X|$. Since $f(X)\subseteq X$ and $|X|$ is finite, $f(X)=X$. Hence $f$ is surjective. Injectivity plus surjectivity equals bijectivity.
\end{proof}

This is sometimes called the ``pigeonhole principle for functions'': on a finite set, you cannot map $N$ distinct inputs to $N$ distinct outputs without covering every output exactly once.

\textbf{Direct corollary:} $f$ is a permutation of $X$. The total number of distinguishable states is strictly conserved under dynamics---not just on average, but exactly at every step.

\subsubsection{A1.2 Cyclic decomposition}

Every permutation $\sigma$ on a finite set of size $M=2^N$ can be decomposed into a product of disjoint cycles:
\begin{equation}
f = \gamma_1 \circ \gamma_2 \circ \cdots \circ \gamma_k,
\end{equation}
where each cycle $\gamma_\ell = (x_\ell^{(1)}, x_\ell^{(2)}, \ldots, x_\ell^{(L_\ell)})$ has length $L_\ell$ and $\sum_{\ell=1}^k L_\ell = 2^N$. The notation $(a_1,a_2,\ldots,a_L)$ means $f(a_i)=a_{i+1}$ for $i<L$ and $f(a_L)=a_1$.

\textbf{Structural consequences:}
\begin{enumerate}[label=(\roman*)]
\item Every state is periodic: $f^{L_\ell}(x) = x$ for every $x$ in cycle $\gamma_\ell$.
\item The dynamics is time-reversible: $f^{-1}$ exists and is also a permutation.
\item No information is ever destroyed---it is only rearranged among the cells.
\item The state space is partitioned into invariant subspaces (cycles), each closed under $f$.
\end{enumerate}

\subsubsection{A1.3 Kac recurrence from cycle structure}

For a subsystem of $N_S$ cells (coarse-graining over the remaining $N_E=N-N_S$ cells), the recurrence time is controlled by the cycle structure of $f$ projected onto the subsystem. The expected recurrence time for a specific subsystem state is:
\begin{equation}
\langle T_{\rm rec}\rangle = 2^{N_S},
\end{equation}
which follows from the uniform distribution over the $2^{N_S}$ possible subsystem macrostates and the fact that $f$ is a permutation (so each macrostate is visited exactly once per full cycle through the relevant part of the decomposition). The independence from $N_E$ is the key result: a large environment accelerates relaxation (the time between recurrence-uncorrelated visits) but does not change the mean recurrence interval.

\textbf{Flag:} The derivation of $\langle T_{\rm rec}\rangle = 2^{N_S}$ assumes that the projection of the permutation onto the subsystem is approximately a random permutation (the ``Kac idealization''). For a specific $f$, the recurrence time can be shorter (if $f$ has short cycles) or longer. The uniform-random-permutation assumption is standard in ergodic theory and is physically justified when the dynamics is sufficiently complex (large $N$, $c$-local interactions with $c>1$). This should be stated explicitly as an approximation.

% ============================================================
\subsection{A2. Derivation of the Continuity Equation $\partial_t q = -\nabla\cdot J$ from Information Conservation}
% ============================================================

\subsubsection{A2.1 Discrete conservation on a graph}

Consider a graph $G=(V,E)$ where each vertex $i\in V$ carries the local vacancy $q_i(t)$. In one discrete time step, information flows along edges. Define $J_{i\to j}^{(t)}$ as the net information transferred from $i$ to $j$ during step $t$ (positive means flow from $i$ to $j$).

The change in $q_i$ over one step is:
\begin{equation}
q_i(t+1) - q_i(t) = -\sum_{j\in\mathcal{N}(i)} J_{i\to j}^{(t)},
\end{equation}
where $\mathcal{N}(i)$ is the set of neighbors of $i$. The minus sign is because outflow from $i$ reduces $q_i$ (makes the cell less vacant). This is an exact accounting identity: information is neither created nor destroyed, only moved.

\subsubsection{A2.2 Continuum limit}

Introduce cell spacing $a$ and time step $\tau$. Define the continuum fields:
\begin{equation}
q(\mathbf{x}_i, t) \equiv q_i(t), \qquad \mathbf{J}(\mathbf{x}_i, t) \equiv \frac{1}{a^{d-1}\tau}\sum_{j} J_{i\to j}^{(t)}\,\hat{\mathbf{e}}_{ij},
\end{equation}
where $\hat{\mathbf{e}}_{ij}$ is the unit vector from $i$ to $j$. The normalization $a^{d-1}\tau$ ensures $\mathbf{J}$ has the dimensions of a flux density.

Taylor-expanding $q_i(t+1) = q(\mathbf{x}_i, t+\tau)$ and $J_{i\to j}^{(t)}$ around $(\mathbf{x}_i, t)$:
\begin{align}
q_i(t+1) - q_i(t) &= \tau\,\partial_t q(\mathbf{x}_i,t) + \mathcal{O}(\tau^2),\\
\sum_{j\in\mathcal{N}(i)} J_{i\to j}^{(t)} &= a^d\,\nabla\cdot\mathbf{J}(\mathbf{x}_i,t) + \mathcal{O}(a^{d+1}),
\end{align}
where the sum over neighbors approximates the divergence by the discrete Gauss theorem. In the limit $a\to 0$, $\tau\to 0$ with $a^2/\tau$ finite:
\begin{equation}
\boxed{\partial_t q = -\nabla\cdot\mathbf{J}}.
\end{equation}

\subsubsection{A2.3 Total $q$ conservation}

Integrating over all space with vanishing flux at infinity:
\begin{equation}
\frac{d}{dt}\int_{\mathbb{R}^d} q(\mathbf{x},t)\,d^d x = -\int_{\mathbb{R}^d} \nabla\cdot\mathbf{J}\,d^d x = -\oint_{S_\infty} \mathbf{J}\cdot d\mathbf{S} = 0,
\end{equation}
provided $\mathbf{J}$ decays faster than $r^{-(d-1)}$ at infinity. Total $q$ is a conserved quantity. This is the continuum expression of the discrete fact that $\sum_i q_i$ changes only by boundary fluxes, which vanish for an isolated system.

\textbf{Approximation flagged:} The continuum limit assumes that $q$ varies slowly on the scale of the cell spacing $a$. Near $q\to 0$ cores, this assumption breaks down and the discrete description is necessary. All continuum conclusions are reliable only in the regime $a|\nabla\ln q|\ll 1$.

% ============================================================
\subsection{A3. Derivation of J-Current Dynamics from Finite Propagation Speed (Cattaneo Model)}
% ============================================================

\subsubsection{A3.1 The infinite-speed problem}

The standard parabolic flux-gradient relation $\mathbf{J} = -\kappa\nabla(\ln q)$ combined with continuity gives $\partial_t q = \kappa\nabla^2(\ln q)$. This is a diffusion equation: an initially localized perturbation instantly affects all space (the fundamental solution is Gaussian with $t>0$ support everywhere). This violates the finite propagation speed $c$ encoded in A3.

\subsubsection{A3.2 Cattaneo's modification of Fourier's law}

Cattaneo (1948) addressed the analogous problem in heat conduction by introducing a relaxation term into Fourier's law:
\begin{equation}
\tau\,\partial_t \mathbf{J} + \mathbf{J} = -\kappa\nabla(\ln q),
\label{eq:cattaneo}
\end{equation}
where $\tau$ is a relaxation time. For $\tau=0$, we recover the original parabolic flux. For $\tau>0$, the flux does not respond instantaneously to the gradient---it relaxes toward the gradient-driven value on timescale $\tau$.

\subsubsection{A3.3 Derivation of the telegraph equation}

Take the time derivative of the continuity equation:
\begin{equation}
\partial_t^2 q = -\nabla\cdot\partial_t\mathbf{J}.
\end{equation}

From the Cattaneo relation (\ref{eq:cattaneo}): $\partial_t\mathbf{J} = -\frac{1}{\tau}\mathbf{J} - \frac{\kappa}{\tau}\nabla(\ln q)$. Substituting:
\begin{equation}
\partial_t^2 q = \nabla\cdot\left(\frac{1}{\tau}\mathbf{J} + \frac{\kappa}{\tau}\nabla(\ln q)\right).
\end{equation}

Using continuity again: $\nabla\cdot\mathbf{J} = -\partial_t q$. Also, $\nabla\cdot(\kappa\nabla(\ln q)/\tau) = (\kappa/\tau)\nabla^2(\ln q)$. Therefore:
\begin{equation}
\boxed{\partial_t^2 q + \frac{1}{\tau}\partial_t q = \frac{\kappa}{\tau}\,\nabla^2(\ln q)}.
\label{eq:telegraph}
\end{equation}

This is the \textbf{telegraph equation} (also called the damped wave equation). It is hyperbolic, not parabolic.

\subsubsection{A3.4 Propagation speed}

Linearizing around a constant background $q_0$: let $q = q_0 + \delta q$ with $|\delta q|\ll q_0$. Then $\ln q \approx \ln q_0 + \delta q/q_0$, so $\nabla^2(\ln q) \approx \nabla^2(\delta q)/q_0$. Equation (\ref{eq:telegraph}) becomes:
\begin{equation}
\partial_t^2(\delta q) + \frac{1}{\tau}\partial_t(\delta q) = c_s^2\,\nabla^2(\delta q),
\end{equation}
where the signal speed is:
\begin{equation}
c_s = \sqrt{\frac{\kappa}{\tau q_0}}.
\end{equation}

This is a standard wave equation with damping. The propagation speed is \textbf{finite}: a localized perturbation at $t=0$ has support only within $r < c_s t$. We identify $c_s$ with $c$ (the speed of light) in the physical correspondence, which relates the microscopics ($\kappa$, $\tau$) to the observed constant $c$:
\begin{equation}
\frac{\kappa}{\tau} = c^2 q_0.
\label{eq:micro_macro}
\end{equation}

\subsubsection{A3.5 Overdamped limit}

When $\tau \ll L^2/\kappa$ (the relaxation time is much shorter than the diffusion time across the system size $L$), the $\partial_t^2 q$ term is negligible. The telegraph equation reduces to the parabolic $q$-field equation. In this limit, propagation appears instantaneous on the timescale of observation, but the underlying dynamics respects the finite speed $c$.

\textbf{The static limit of the telegraph equation is identical to the static limit of the parabolic equation:} setting $\partial_t = \partial_t^2 = 0$ in (\ref{eq:telegraph}) gives $\nabla^2(\ln q)=0$, recovering Eq.~(A5). The gravitational correspondence is therefore insensitive to whether we use the parabolic or hyperbolic form.

\subsubsection{A3.6 Cattaneo number}

Define the Cattaneo number $\mathcal{C} \equiv \tau\kappa/L^2q_0 = (c_s\tau/L)^2$. The parabolic approximation is valid when $\mathcal{C} \ll 1$. In the opposite regime $\mathcal{C} \gg 1$, wave-like propagation dominates. For the $q$-field, with $c_s=c$, $\tau \sim \ell_P/c$ (Planck time), and $L$ as the system size:
\begin{equation}
\mathcal{C} \sim \left(\frac{\ell_P}{L}\right)^2 \sim 10^{-70} \quad (\text{for laboratory scales}),
\end{equation}
so the parabolic approximation is excellent except at the Planck scale.

\textbf{Flag:} The identification $\tau \sim \ell_P/c$ is a dimensional estimate. The actual value of $\tau$ is not derived within the framework. Equation (\ref{eq:micro_macro}) provides one relation between $\kappa$ and $\tau$, leaving one degree of freedom.

% ============================================================
\subsection{A4. Full Expansion of $\nabla^2(\ln q) = \nabla^2 q/q - |\nabla q|^2/q^2$}
% ============================================================

\subsubsection{A4.1 Vector identity derivation}

Let $\psi \equiv \ln q$, so $q = e^\psi$. By the chain rule:
\begin{equation}
\nabla q = \nabla(e^\psi) = e^\psi \nabla\psi = q\,\nabla\psi.
\end{equation}

Therefore:
\begin{equation}
\nabla\psi = \frac{\nabla q}{q}.
\end{equation}

Take the divergence:
\begin{align}
\nabla^2\psi = \nabla\cdot(\nabla\psi) &= \nabla\cdot\left(\frac{\nabla q}{q}\right)\\
&= \frac{\nabla\cdot(\nabla q)}{q} + (\nabla q)\cdot\nabla\left(\frac{1}{q}\right)\\
&= \frac{\nabla^2 q}{q} + (\nabla q)\cdot\left(-\frac{\nabla q}{q^2}\right)\\
&= \frac{\nabla^2 q}{q} - \frac{|\nabla q|^2}{q^2}.
\end{align}

Hence:
\begin{equation}
\boxed{\nabla^2(\ln q) = \frac{\nabla^2 q}{q} - \frac{|\nabla q|^2}{q^2}}.
\label{eq:laplacian_expansion}
\end{equation}

\subsubsection{A4.2 Alternative form}

An equivalent and sometimes more convenient expression is:
\begin{equation}
\nabla^2(\ln q) = \nabla\cdot\left(\frac{\nabla q}{q}\right) = \frac{1}{q}\nabla^2 q - \frac{1}{q^2}|\nabla q|^2.
\end{equation}

The $q$-field equation can therefore be written in three equivalent forms:
\begin{align}
\partial_t q &= \kappa\,\nabla^2(\ln q) &&\text{(compact form)},\\
\partial_t q &= \kappa\left(\frac{\nabla^2 q}{q} - \frac{|\nabla q|^2}{q^2}\right) &&\text{(expanded form)},\\
\partial_t q &= \kappa\,\nabla\cdot\left(\frac{\nabla q}{q}\right) &&\text{(divergence form)}.
\end{align}

\subsubsection{A4.3 The $1/q$ amplification mechanism}

The expanded form reveals the physical mechanism:
\begin{itemize}
\item The term $\nabla^2 q/q$ is the standard diffusion term, but amplified by $1/q$: in low-$q$ regions, even small Laplacians produce large effects.
\item The term $-|\nabla q|^2/q^2$ is a nonlinear drift/advection term: it drives $q$ down gradients (toward regions of lower $q$), amplifying inhomogeneity---this is anti-diffusion.
\end{itemize}

Consider a Gaussian bump $q(r) = q_0 + \delta q\,e^{-r^2/2\sigma^2}$. Then:
\begin{align}
\nabla^2 q &\sim -\frac{\delta q}{\sigma^2}, &
|\nabla q|^2 &\sim \frac{(\delta q)^2 r^2}{\sigma^4}e^{-r^2/\sigma^2}.
\end{align}

The ratio of the anti-diffusive to diffusive terms scales as $(\delta q/q)(r^2/\sigma^2)\cdot(\delta q/q)$. When $q\ll 1$, this ratio is large: the anti-diffusion dominates. This mathematically captures the `'amplification toward $q\to 0$'' discussed in the main text.

\textbf{No approximations in A4.} Equation (\ref{eq:laplacian_expansion}) is an exact vector calculus identity valid wherever $q>0$ and $q$ is twice differentiable.

% ============================================================
\subsection{A5. Spherical Static Solution: From $\nabla^2(\ln 1/q)=0$ to $q(r)=e^{-GM/rc^2}$}
% ============================================================

\subsubsection{A5.1 Static equation in terms of $\psi$}

Define the information potential $\psi \equiv -\ln q$. Since $\ln(1/q) = -\ln q = \psi$, the static equation $\nabla^2(\ln 1/q)=0$ is:
\begin{equation}
\nabla^2\psi = 0.
\label{eq:static_psi}
\end{equation}

This is the Laplace equation for $\psi$.

\subsubsection{A5.2 Spherical symmetry}

Assume $\psi = \psi(r)$, $r=|\mathbf{x}|$. The Laplace operator in spherical coordinates:
\begin{equation}
\nabla^2\psi = \frac{1}{r^2}\frac{d}{dr}\left(r^2\frac{d\psi}{dr}\right) = 0.
\end{equation}

Multiply by $r^2$ and integrate once:
\begin{equation}
r^2\frac{d\psi}{dr} = C_1 = \text{const}.
\label{eq:first_integral}
\end{equation}

\subsubsection{A5.3 Determining $C_1$ via Gauss's law with source}

To introduce a source, we generalize the static equation to include a source term representing locked information (mass):
\begin{equation}
\nabla^2\psi = -\frac{4\pi G}{c^2}\,\rho_m(\mathbf{x}),
\label{eq:poisson_psi}
\end{equation}
where $\rho_m$ is the mass density and the coupling constant is calibrated to match Newtonian gravity (see \S A6). For a point mass $M$ at the origin, $\rho_m(\mathbf{x}) = M\delta^{(3)}(\mathbf{x})$.

Integrate (\ref{eq:poisson_psi}) over a sphere of radius $r$:
\begin{equation}
\int_{B_r} \nabla^2\psi\,d^3x = -\frac{4\pi G}{c^2}M.
\end{equation}

By the divergence theorem:
\begin{equation}
\oint_{S_r} \nabla\psi\cdot d\mathbf{S} = 4\pi r^2\frac{d\psi}{dr} = -\frac{4\pi G}{c^2}M.
\end{equation}

Therefore:
\begin{equation}
\frac{d\psi}{dr} = -\frac{GM}{c^2 r^2}, \qquad C_1 = -\frac{GM}{c^2}.
\label{eq:C1_value}
\end{equation}

\textbf{Key point:} Equation (\ref{eq:poisson_psi}) is the \emph{definition} of how mass sources the $\psi$-field. The $1/r^2$ force law follows, but the proportionality constant $G$ is calibrated, not derived.

\subsubsection{A5.4 Second integration}

From $d\psi/dr = -GM/(c^2 r^2)$:
\begin{equation}
\psi(r) = \frac{GM}{c^2 r} + C_2.
\end{equation}

\subsubsection{A5.5 Boundary condition: $q(\infty)=1$}

At spatial infinity, the influence of any finite mass vanishes: the field should approach the vacuum state $q=1$. Therefore $\psi(\infty) = -\ln(1) = 0$. This forces $C_2 = 0$:
\begin{equation}
\psi(r) = \frac{GM}{c^2 r}.
\end{equation}

\subsubsection{A5.6 Recovering $q(r)$}

Since $\psi = -\ln q$, we have $q = e^{-\psi}$:
\begin{equation}
\boxed{q(r) = e^{-GM/rc^2}}.
\label{eq:q_spherical}
\end{equation}

\subsubsection{A5.7 Properties of the solution}

\begin{enumerate}[label=(\roman*)]
\item $q(r) \to 1$ as $r\to\infty$ (asymptotic flatness).
\item $q(r) \to 0$ as $r\to 0$ (information capacity exhaustion at the center).
\item $q(r) > 0$ for all $r > 0$ (no sharp event horizon---$q$ never reaches exactly zero at any finite $r$).
\item The solution is self-similar: $q(\lambda r; \lambda M) = q(r; M)$ for any $\lambda > 0$ (the exponent depends only on $M/r$).
\end{enumerate}

\subsubsection{A5.8 The Newtonian correspondence via the $1/q$ potential}

For completeness, note that if one uses $\varphi(q)=1/q$ instead of $\varphi(q)=-\ln q$, the static equation $\nabla^2(1/q)=0$ yields the spherical solution:
\begin{equation}
\frac{1}{q(r)} = 1 + \frac{GM}{c^2 r}, \qquad q(r) = \frac{1}{1 + GM/rc^2}.
\end{equation}

The exponential form (\ref{eq:q_spherical}) is the solution for $\varphi(q)=-\ln q$; the rational form is the solution for $\varphi(q)=1/q$. Both belong to the same universality class $\Phi$ and produce identical qualitative physics. The difference becomes quantitative only at strong fields, where neither form is reliable (the discrete treatment is required).

\textbf{Approximation flagged:} The derivation assumes a \textbf{point mass} source. For extended mass distributions, one solves $\nabla^2\psi = -(4\pi G/c^2)\rho_m(\mathbf{x})$ with $\rho_m$ extended, and the solution is the convolution of the point-mass Green's function with $\rho_m$:
\begin{equation}
\psi(\mathbf{x}) = \frac{G}{c^2}\int \frac{\rho_m(\mathbf{x}')}{|\mathbf{x}-\mathbf{x}'|}\,d^3x'.
\end{equation}

% ============================================================
\subsection{A6. Weak-Field Expansion: $q(r)\approx 1-GM/rc^2$ $\to$ Newtonian Limit}
% ============================================================

\subsubsection{A6.1 Taylor expansion}

For $GM/rc^2 \ll 1$ (weak field, large $r$), expand the exponential:
\begin{equation}
q(r) = e^{-GM/rc^2} = 1 - \frac{GM}{rc^2} + \frac{1}{2}\left(\frac{GM}{rc^2}\right)^2 - \cdots
\end{equation}

To first order:
\begin{equation}
\boxed{q(r) \approx 1 - \frac{GM}{rc^2}}.
\label{eq:weak_field}
\end{equation}

\subsubsection{A6.2 Newtonian gravitational potential}

The Newtonian gravitational potential for a point mass is:
\begin{equation}
\Phi_N(r) = -\frac{GM}{r}.
\end{equation}

Comparing with (\ref{eq:weak_field}):
\begin{equation}
q(r) \approx 1 + \frac{\Phi_N(r)}{c^2}.
\end{equation}

Thus, to first order in the weak-field expansion:
\begin{equation}
\boxed{\frac{\Phi_N}{c^2} = q - 1 = -(1-q)}.
\label{eq:newton_correspondence}
\end{equation}

The Newtonian potential (in units of $c^2$) equals \emph{minus the deviation of $q$ from unity}---i.e., minus the occupied fraction $1-q$.

\subsubsection{A6.3 Gravitational acceleration}

The Newtonian gravitational field (acceleration) is:
\begin{equation}
\mathbf{g} = -\nabla\Phi_N = -c^2\nabla q.
\end{equation}

For the spherical solution: $\mathbf{g} = -c^2\,d(e^{-GM/rc^2})/dr\,\hat{\mathbf{r}} = -(GM/r^2)e^{-GM/rc^2}\,\hat{\mathbf{r}} \approx -(GM/r^2)\,\hat{\mathbf{r}}$ in the weak-field limit. This recovers the inverse-square law.

\subsubsection{A6.4 Poisson equation recovery}

From $\Phi_N = c^2(q-1)$ and the static $q$-field equation $\nabla^2(\ln q)=0$:

Expand $\ln q = \ln(1 + \Phi_N/c^2) \approx \Phi_N/c^2$ for $|\Phi_N|/c^2 \ll 1$. Then $\nabla^2(\ln q) \approx \nabla^2(\Phi_N)/c^2 = 0$, so $\nabla^2\Phi_N=0$ in vacuum. With sources (as derived in \S A5.3):
\begin{equation}
\nabla^2\Phi_N = 4\pi G\rho_m,
\end{equation}
which is exactly the Newtonian Poisson equation.

\subsubsection{A6.5 Validity regime}

The weak-field approximation requires:
\begin{equation}
\frac{GM}{rc^2} \ll 1 \quad\Longleftrightarrow\quad r \gg \frac{GM}{c^2} \equiv \frac{r_s}{2},
\end{equation}
where $r_s = 2GM/c^2$ is the Schwarzschild radius. For the Earth, $GM/rc^2 \sim 7\times 10^{-10}$ at the surface---the weak-field approximation is excellent. For a neutron star, $GM/Rc^2 \sim 0.21$---the approximation is marginal. For a black hole at the horizon, $GM/r_s c^2 = 1/2$---the approximation is unreliable.

\textbf{Approximation flagged:} The identification $\Phi_N = c^2(q-1)$ in (\ref{eq:newton_correspondence}) is valid only to first order in $GM/rc^2$. The next-to-leading-order term differs between the exponential and rational forms, and neither can be trusted without a full strong-field treatment of the discrete dynamics.

% ============================================================
\subsection{A7. Energy Functional $E[q,J]$ and Dissipation Theorem $dE/dt \le 0$}
% ============================================================

\subsubsection{A7.1 Definition of the free information functional}

Define the Lyapunov functional (``free information''):
\begin{equation}
\mathcal{F}[q] \equiv \int_{\mathbb{R}^d} \Phi(q(\mathbf{x}))\,d^d x,
\end{equation}
where $\Phi'(q) = \ln q$, i.e.:
\begin{equation}
\Phi(q) = \int^q \ln s\,ds = q\ln q - q + \text{const}.
\end{equation}

We fix the constant by requiring $\Phi(1)=0$ (vacuum has zero free information):
\begin{equation}
\Phi(q) = q\ln q - q + 1.
\end{equation}

Note that $\Phi(q) \ge 0$ for $q\in[0,1]$, with $\Phi(0)=1$ and $\Phi(1)=0$.

\subsubsection{A7.2 Time derivative}

\begin{align}
\frac{d\mathcal{F}}{dt} &= \int \Phi'(q)\,\partial_t q\,d^d x\\
&= \int (\ln q)\,\partial_t q\,d^d x.
\end{align}

Using the $q$-field equation $\partial_t q = \kappa\nabla^2(\ln q)$:
\begin{align}
\frac{d\mathcal{F}}{dt} &= \kappa\int (\ln q)\,\nabla^2(\ln q)\,d^d x\\
&= \kappa\int \left[\nabla\cdot((\ln q)\nabla(\ln q)) - |\nabla(\ln q)|^2\right] d^d x\\
&= \kappa\oint_{S_\infty} (\ln q)\nabla(\ln q)\cdot d\mathbf{S} - \kappa\int |\nabla(\ln q)|^2\,d^d x,
\end{align}
where we used integration by parts (Green's first identity).

The boundary term vanishes for isolated systems ($q\to 1$, $\nabla(\ln q)\to 0$ at infinity). Therefore:
\begin{equation}
\boxed{\frac{d\mathcal{F}}{dt} = -\kappa\int_{\mathbb{R}^d} |\nabla(\ln q)|^2\,d^d x \le 0}.
\label{eq:dissipation}
\end{equation}

\subsubsection{A7.3 Physical interpretation}

\begin{enumerate}[label=(\roman*)]
\item $\mathcal{F}[q]$ is a monotonically non-increasing functional of the $q$-field. The system evolves irreversibly toward configurations of lower free information.
\item Dissipation ceases ($d\mathcal{F}/dt = 0$) if and only if $\nabla(\ln q)=0$ everywhere---i.e., $q$ is spatially uniform. The static equilibrium is the state of uniform $q$.
\item The rate of dissipation is proportional to the integrated squared gradient of the information potential. Steep gradients dissipate rapidly.
\item This is an $H$-theorem for the $q$-field: the dynamics has a built-in arrow of time.
\end{enumerate}

\subsubsection{A7.4 Generalization to arbitrary flux $\varphi \in \Phi$}

For a general flux potential $\varphi(q)$ with $\varphi'(q)<0$, the PDE is $\partial_t q = -\nabla^2\varphi(q)$. Define $\Phi_\varphi(q) = \int^q \varphi(s)\,ds$. Then:
\begin{equation}
\frac{d}{dt}\int \Phi_\varphi(q)\,d^d x = -\int |\nabla\varphi(q)|^2\,d^d x \le 0,
\end{equation}
so every $\varphi\in\Phi$ admits a Lyapunov functional. The dissipation is universal.

\subsubsection{A7.5 The energy functional $E[q,\mathbf{J}]$ for the Cattaneo extension}

For the hyperbolic extension $\tau\partial_t\mathbf{J} + \mathbf{J} = -\kappa\nabla(\ln q)$ with $\partial_t q = -\nabla\cdot\mathbf{J}$, define:
\begin{equation}
E[q,\mathbf{J}] \equiv \int\left[\Phi(q) + \frac{\tau}{2\kappa}|\mathbf{J}|^2\right] d^d x,
\end{equation}
where the second term is the ``kinetic energy'' of the information current.

The time derivative:
\begin{align}
\frac{dE}{dt} &= \int\left[\Phi'(q)\partial_t q + \frac{\tau}{\kappa}\mathbf{J}\cdot\partial_t\mathbf{J}\right] d^d x\\
&= \int\left[-\ln q\,\nabla\cdot\mathbf{J} + \frac{\tau}{\kappa}\mathbf{J}\cdot\left(-\frac{1}{\tau}\mathbf{J} - \frac{\kappa}{\tau}\nabla(\ln q)\right)\right] d^d x\\
&= \int\left[-\ln q\,\nabla\cdot\mathbf{J} - \frac{1}{\kappa}|\mathbf{J}|^2 - \mathbf{J}\cdot\nabla(\ln q)\right] d^d x.
\end{align}

The first and third terms combine: $-\ln q\,\nabla\cdot\mathbf{J} - \mathbf{J}\cdot\nabla(\ln q) = -\nabla\cdot(\ln q\,\mathbf{J})$, a total divergence that vanishes on integration. Hence:
\begin{equation}
\boxed{\frac{dE}{dt} = -\frac{1}{\kappa}\int |\mathbf{J}|^2\,d^d x \le 0}.
\label{eq:energy_dissipation}
\end{equation}

The energy $E[q,\mathbf{J}]$ is a strict Lyapunov functional for the Cattaneo system: dissipation is proportional to the total squared current, and equilibrium is reached only when $\mathbf{J}=0$ everywhere.

\textbf{No approximations in A7.} Equations (\ref{eq:dissipation}) and (\ref{eq:energy_dissipation}) are exact consequences of the respective PDE systems, assuming sufficient regularity for integration by parts.

%%%%%%%%%%%%%%%%%%%%%%%%%%%%%%%%%%%%%%%%%%%%%%%%%%%%%%%%%%%%%%%%%%%%%
\section{Mathematical Rigor}
%%%%%%%%%%%%%%%%%%%%%%%%%%%%%%%%%%%%%%%%%%%%%%%%%%%%%%%%%%%%%%%%%%%%%

% ============================================================
\subsection{B1. Well-Posedness Proof for the Coupled $(q,\mathbf{J})$ System}
% ============================================================

\subsubsection{B1.1 The system}

We consider the hyperbolic (Cattaneo) system on $\mathbb{R}^d$:
\begin{align}
\partial_t q &= -\nabla\cdot\mathbf{J}, \label{eq:sys1}\\
\tau\partial_t\mathbf{J} + \mathbf{J} &= -\kappa\nabla(\ln q), \label{eq:sys2}
\end{align}
with initial conditions $q(\mathbf{x},0)=q_0(\mathbf{x})$, $\mathbf{J}(\mathbf{x},0)=\mathbf{J}_0(\mathbf{x})$.

\subsubsection{B1.2 Local well-posedness}

\textbf{Theorem B1.} Let $q_0 \in H^s(\mathbb{R}^d)$ with $q_0(\mathbf{x}) \ge q_{\min} > 0$ for all $\mathbf{x}$, and $\mathbf{J}_0 \in H^{s-1}(\mathbb{R}^d)$ with $s > d/2+1$. Then there exists $T>0$ such that the system (\ref{eq:sys1})-(\ref{eq:sys2}) has a unique solution $(q,\mathbf{J})$ with:
\begin{equation}
q \in C([0,T]; H^s(\mathbb{R}^d)) \cap C^1([0,T]; H^{s-2}(\mathbb{R}^d)), \quad \mathbf{J} \in C([0,T]; H^{s-1}(\mathbb{R}^d)).
\end{equation}

\textbf{Proof sketch.} Rewrite (\ref{eq:sys1})-(\ref{eq:sys2}) as a first-order symmetric hyperbolic system in the variables $(q, \mathbf{J})$. The nonlinearity $\nabla(\ln q) = \nabla q/q$ is smooth and bounded for $q \ge q_{\min} > 0$. Standard theory (Kato, Majda) for symmetric hyperbolic systems with smooth coefficients then guarantees local existence and uniqueness.

The strict positivity condition $q_0 \ge q_{\min} > 0$ is essential: as $q\to 0$, the nonlinearity $\nabla(\ln q) = \nabla q/q$ loses Lipschitz continuity, and the proof fails. This is the mathematical reflection of the physical fact that the $q\to 0$ singularity requires a discrete treatment.

\subsubsection{B1.3 Global existence}

\textbf{Theorem B1 (continued).} If the initial data further satisfies $q_0 \ge \delta > 0$ with a uniform positive lower bound, then the solution exists for all $t > 0$ and satisfies the energy estimate:
\begin{equation}
E(t) + \frac{1}{\kappa}\int_0^t \int |\mathbf{J}|^2\,d^d x\,dt' = E(0).
\end{equation}

\textbf{Proof sketch.} The energy $E(t)$ (defined in \S A7.5) is non-increasing. Since $\Phi(q) \ge 0$ and $|\mathbf{J}|^2 \ge 0$, both $q$ and $\mathbf{J}$ remain bounded in the energy norm. A continuation argument extends the local solution to global.

\subsubsection{B1.4 The $q\to 0$ problem}

When $\inf q_0 = 0$, the hyperbolic system is not locally well-posed in Sobolev spaces. The formation of $q=0$ singularities corresponds to the development of ``classical cores''---regions where the continuum description breaks down. In the discrete DGF model, $q=0$ is well-defined (all cells occupied), and the dynamics continues. The continuum PDE should be understood as an approximation valid away from $q=0$ cores.

\textbf{Flag:} Rigorous well-posedness is proved only for $q \ge q_{\min} > 0$. The formation and dynamics of $q=0$ cores require the discrete theory.

% ============================================================
\subsection{B2. Spectral Analysis: Eigenvalues, Stability, Dispersion Relation}
% ============================================================

\subsubsection{B2.1 Linearization}

Linearize $\partial_t q = \kappa\nabla^2(\ln q)$ around a constant background $q_0$: let $q = q_0 + \epsilon\eta$ with $\epsilon\ll 1$. Then:
\begin{equation}
\ln q = \ln q_0 + \ln(1 + \epsilon\eta/q_0) \approx \ln q_0 + \frac{\epsilon\eta}{q_0},
\end{equation}
so $\nabla^2(\ln q) \approx (\epsilon/q_0)\nabla^2\eta$. The linearized equation is:
\begin{equation}
\partial_t\eta = \frac{\kappa}{q_0}\nabla^2\eta.
\label{eq:linearized}
\end{equation}

This is the heat equation with diffusivity $\kappa/q_0 = D(q_0)$. Consistent: $D(q_0)\to\infty$ as $q_0\to 0$ (fast diffusion), $D(1)=\kappa$ (finite vacuum diffusion).

\subsubsection{B2.2 Fourier modes and dispersion}

Seek solutions $\eta(\mathbf{x},t) = \hat{\eta}_{\mathbf{k}}\,e^{i\mathbf{k}\cdot\mathbf{x} + \sigma t}$. Substituting into (\ref{eq:linearized}):
\begin{equation}
\sigma(\mathbf{k}) = -\frac{\kappa}{q_0}|\mathbf{k}|^2.
\end{equation}

All modes with $\mathbf{k}\neq 0$ decay exponentially ($\sigma < 0$). The $\mathbf{k}=0$ mode is marginal ($\sigma=0$), corresponding to conservation of total $q$. The system is linearly stable.

\subsubsection{B2.3 Hyperbolic dispersion}

For the Cattaneo system, seek solutions with $\partial_t \to \sigma$, $\nabla \to i\mathbf{k}$. From (\ref{eq:telegraph}):
\begin{equation}
\sigma^2 + \frac{1}{\tau}\sigma + \frac{\kappa}{\tau q_0}|\mathbf{k}|^2 = 0.
\end{equation}

Solving:
\begin{equation}
\sigma_\pm(\mathbf{k}) = -\frac{1}{2\tau} \pm \sqrt{\frac{1}{4\tau^2} - \frac{\kappa}{\tau q_0}|\mathbf{k}|^2}.
\end{equation}

\textbf{Dispersion regimes:}
\begin{enumerate}[label=(\roman*)]
\item \textbf{Long wavelength} ($|\mathbf{k}| < k_{\rm cross}$): Both roots are real and negative---purely decaying (overdamped) modes. Where $k_{\rm cross} \equiv (4\kappa\tau/q_0)^{-1/2} = (2c_s\tau)^{-1}$.
\item \textbf{Short wavelength} ($|\mathbf{k}| > k_{\rm cross}$): Complex conjugate roots with $\Re(\sigma) = -1/(2\tau)$---damped oscillatory modes. Phase velocity: $v_\phi = \Im(\sigma)/|\mathbf{k}| \to c_s$ as $|\mathbf{k}|\to\infty$. Group velocity: $v_g = d\Im(\sigma)/d|\mathbf{k}| \to c_s$ as $|\mathbf{k}|\to\infty$.
\end{enumerate}

The crossover wavenumber $k_{\rm cross}$ separates diffusive (long-wavelength) from wave-like (short-wavelength) behavior. For $|\mathbf{k}|\ll k_{\rm cross}$, expand:
\begin{equation}
\sigma_+ \approx -\frac{\kappa}{q_0}|\mathbf{k}|^2, \qquad \sigma_- \approx -\frac{1}{\tau},
\end{equation}
recovering the parabolic dispersion plus a rapidly decaying mode.

\subsubsection{B2.4 Eigenvalue problem for the static equation}

In a bounded domain $\Omega$ with boundary conditions $q|_{\partial\Omega} = q_0$ (or Neumann $\partial_n q|_{\partial\Omega}=0$), the static equation $\nabla^2(\ln q) = 0$ with $\ln q \equiv \psi$ is the Dirichlet (or Neumann) eigenvalue problem for the Laplacian. The eigenvalues $\lambda_n$ of $-\nabla^2$ on $\Omega$ satisfy $0 = \lambda_0 < \lambda_1 \le \lambda_2 \le \cdots$. The static solution exists uniquely for Dirichlet conditions; for Neumann conditions, it exists up to an additive constant in $\psi$, which is fixed by the total $q$ conservation.

\textbf{No approximations in B2.} The dispersion relations are exact for the linearized equations. Relevance to the full nonlinear system is restricted to the linear regime $|\delta q| \ll q_0$.

% ============================================================
\subsection{B3. Overdamped Limit as Singular Perturbation ($\gamma\to\infty$, $c^2/\gamma$ Fixed)}
% ============================================================

\subsubsection{B3.1 Scaling}

Rewrite the telegraph equation (\ref{eq:telegraph}) in terms of the damping rate $\gamma \equiv 1/\tau$ and the squared propagation speed $c_s^2 \equiv \kappa/(\tau q_0) = \gamma\kappa/q_0$:
\begin{equation}
\frac{1}{\gamma}\,\partial_t^2 q + \partial_t q = \frac{c_s^2}{\gamma}\,\nabla^2(\ln q).
\end{equation}

\textbf{The overdamped limit:} $\gamma \to \infty$ with $c_s^2/\gamma = \kappa/q_0$ held fixed. In this limit, the inertial term $\gamma^{-1}\partial_t^2 q$ is negligible, and we recover:
\begin{equation}
\partial_t q = \frac{c_s^2}{\gamma}\nabla^2(\ln q) = \frac{\kappa}{q_0}\nabla^2(\ln q),
\end{equation}
which is the parabolic $q$-field equation.

\subsubsection{B3.2 Singular perturbation structure}

This is a singular perturbation: the order of the PDE drops from 2 (hyperbolic, requires two initial conditions $q_0$ and $\partial_t q_0$) to 1 (parabolic, requires only $q_0$). The second initial condition is lost in a boundary layer of thickness $\mathcal{O}(1/\gamma)$ at $t=0$.

Formally, the solution has the asymptotic expansion:
\begin{equation}
q(\mathbf{x},t;\gamma) = q^{(0)}(\mathbf{x},t) + e^{-\gamma t}q^{(1)}(\mathbf{x},\gamma t) + \mathcal{O}(1/\gamma),
\end{equation}
where $q^{(0)}$ solves the parabolic equation and the initial layer $q^{(1)}$ decays exponentially on the fast timescale $\tau = 1/\gamma$.

\subsubsection{B3.3 Matching condition}

For the expansion to be uniform, the initial condition for the outer (parabolic) solution must satisfy:
\begin{equation}
q^{(0)}(\mathbf{x},0) = q_0(\mathbf{x}),
\end{equation}
while the initial layer corrects the initial velocity $\partial_t q(\mathbf{x},0)$ to be consistent with the reduced equation. This is standard boundary-layer theory.

\textbf{Physical interpretation:} On timescales $t \gg \tau$, the system ``forgets'' the initial current configuration $\mathbf{J}_0$ and follows the gradient-driven parabolic dynamics. Since $\tau \sim \ell_P/c \sim 10^{-43}$~s for Planck-scale parameters, the overdamped limit is an extraordinarily good approximation for all astrophysical and laboratory timescales.

% ============================================================
\subsection{B4. Continuum Limit of Graph Laplacian on Locally Finite Graphs}
% ============================================================

\subsubsection{B4.1 Discrete setup}

Consider an infinite regular grid $\Lambda \subset \mathbb{Z}^d$ with spacing $a$. The graph Laplacian acting on a function $u_i = u(\mathbf{x}_i)$ is:
\begin{equation}
(\mathcal{L}_a u)_i \equiv \frac{1}{a^2}\sum_{j \sim i} (u_j - u_i),
\end{equation}
where $j\sim i$ denotes nearest neighbors. The normalization $1/a^2$ ensures the correct continuum limit.

\subsubsection{B4.2 Taylor expansion}

For a smooth function $u(\mathbf{x})$, Taylor-expand around $\mathbf{x}_i$:
\begin{equation}
u(\mathbf{x}_i \pm a\hat{\mathbf{e}}_\alpha) = u(\mathbf{x}_i) \pm a\partial_\alpha u(\mathbf{x}_i) + \frac{a^2}{2}\partial_\alpha^2 u(\mathbf{x}_i) \pm \frac{a^3}{6}\partial_\alpha^3 u(\mathbf{x}_i) + \mathcal{O}(a^4).
\end{equation}

The graph Laplacian (summing over $\alpha = 1,\ldots,d$, with two neighbors per direction):
\begin{align}
(\mathcal{L}_a u)_i &= \frac{1}{a^2}\sum_{\alpha=1}^d \Big[(u(\mathbf{x}_i + a\hat{\mathbf{e}}_\alpha) - u(\mathbf{x}_i)) + (u(\mathbf{x}_i - a\hat{\mathbf{e}}_\alpha) - u(\mathbf{x}_i))\Big]\\
&= \frac{1}{a^2}\sum_{\alpha=1}^d \left[a^2\partial_\alpha^2 u + \mathcal{O}(a^4)\right]\\
&= \nabla^2 u(\mathbf{x}_i) + \mathcal{O}(a^2).
\label{eq:graph_laplacian_limit}
\end{align}

The $\mathcal{O}(a^3)$ terms cancel by symmetry (central difference). The leading error is $\mathcal{O}(a^2)$, proportional to $\sum_\alpha \partial_\alpha^4 u$.

\subsubsection{B4.3 Locally finite graphs}

A graph is \textbf{locally finite} if every vertex has finite degree. For the $q$-field model, the degree is $2d$ (regular grid). The result (\ref{eq:graph_laplacian_limit}) generalizes: for any locally finite graph with a periodic structure, the discrete Laplacian converges to the continuum Laplacian as the lattice spacing tends to zero, provided the graph is ``sufficiently uniform'' (bounded geometry, polynomial volume growth).

For irregular graphs (e.g., random geometric graphs), the convergence is to a weighted Laplacian $\nabla\cdot(\sigma(\mathbf{x})\nabla)$ where $\sigma$ encodes local connectivity variations. The DGF framework assumes a regular grid for simplicity; irregularity would introduce effective permittivity variations in the gravitational analogy.

\subsubsection{B4.4 The $q$-field PDE as a gradient flow on graphs}

The discrete $q$-field dynamics is:
\begin{equation}
\partial_t q_i = \kappa\sum_{j\sim i} (\ln q_j - \ln q_i) = -\kappa(\mathcal{L}\ln\mathbf{q})_i,
\end{equation}
where $(\ln\mathbf{q})_i \equiv \ln q_i$. In the continuum limit $a\to 0$, this becomes:
\begin{equation}
\partial_t q = -\kappa a^2\mathcal{L}_a(\ln q) \to -\kappa\nabla^2(\ln q) = \kappa\nabla^2(\ln q),
\end{equation}
with $\kappa$ absorbing the $a^2$ factor from the graph Laplacian normalization.

\textbf{Approximation flagged:} The convergence $\mathcal{L}_a \to \nabla^2$ is pointwise in $C^4$ functions. Near $q=0$ cores where $q$ varies rapidly, the higher-order terms $\mathcal{O}(a^2\nabla^4\ln q)$ may be significant.

% ============================================================
\subsection{B5. Discrete-to-Continuum Error Estimates ($N\sim 10^{90}$ Cells, Error $\mathcal{O}(10^{-90})$)}
% ============================================================

\subsubsection{B5.1 Cell counting}

The observable universe has volume $V \sim (4\pi/3)(c/H_0)^3 \sim 10^{80}$ m$^3$. If the fundamental cell has Planck volume $\ell_P^3 \sim (1.6\times 10^{-35}~\text{m})^3 \sim 4\times 10^{-105}$ m$^3$, then:
\begin{equation}
N = \frac{V}{\ell_P^3} \sim \frac{10^{80}}{4\times 10^{-105}} \sim 2.5\times 10^{184}.
\label{eq:N_estimate}
\end{equation}

The more conservative estimate $N\sim 10^{90}$ corresponds to a cell size of $\sim 10^{-15}$~m (nuclear scale), which is the scale where quantum field theory degrees of freedom become relevant. The exact cell scale is not determined within DGF; $N$ is a free parameter. We adopt $N\sim 10^{90}$ as a fiducial value for error estimates.

\subsubsection{B5.2 Discretization error scaling}

For a smooth function $u$ on a grid of $N$ cells (spacing $a \sim N^{-1/d}$), the discrete Laplacian approximates the continuum Laplacian with error (from \S B4):
\begin{equation}
\|(\mathcal{L}_a - \nabla^2)u\|_\infty \le C a^2 \|\nabla^4 u\|_\infty \sim \mathcal{O}(N^{-2/d}).
\end{equation}

In $d=3$: error $\sim \mathcal{O}(N^{-2/3})$. For $N\sim 10^{90}$: error $\sim 10^{-60}$. This is the pointwise error per cell.

\subsubsection{B5.3 Global error estimate}

For the static spherical solution $q(r) = e^{-GM/rc^2}$, the fourth derivative scales as:
\begin{equation}
\nabla^4 q \sim \frac{GM}{c^2 r^5} \sim \frac{r_s}{r^5},
\end{equation}
where $r_s$ is the Schwarzschild radius. The error is largest at small $r$. The continuum approximation is reliable when:
\begin{equation}
a^2 \frac{r_s}{r^5} \ll 1 \quad\Longleftrightarrow\quad r \gg (r_s a^2)^{1/5}.
\end{equation}

For a solar-mass black hole ($r_s \sim 3\times 10^3$~m) and Planck-scale cells ($a \sim 10^{-35}$~m): $(r_s a^2)^{1/5} \sim (3\times 10^3 \times 10^{-70})^{1/5} \sim 10^{-13}$~m, far below any physically accessible scale. The continuum approximation is excellent down to sub-nuclear distances.

\subsubsection{B5.4 Finite-size effects}

For a finite system of $N$ cells, the total number of microstates is $2^N$. The entropy $S = N\ln 2$. Fluctuations in macroscopic observables (like $q$) scale as $1/\sqrt{N}$. For $N\sim 10^{90}$:
\begin{equation}
\frac{\Delta q}{q} \sim \frac{1}{\sqrt{N}} \sim 10^{-45}.
\end{equation}

Finite-size fluctuations are astronomically suppressed. The continuum deterministic PDE is an extraordinarily accurate description of the expected $q$-field. Fluctuations matter only when probing at the level of $10^{45}$ standard deviations---a practical impossibility.

\subsubsection{B5.5 Important caveat}

\textbf{The error estimate $\mathcal{O}(10^{-90})$ refers to the finite-$N$ correction to the continuum PDE, not to the absolute accuracy of the DGF framework.} If the DGF framework itself is an approximation to a deeper theory, this error estimate does not capture that mismatch. The $\mathcal{O}(10^{-90})$ estimate says: \emph{if} the discrete DGF model on $N=10^{90}$ cells is the exact microscopic theory, \emph{then} the continuum PDE approximates it with error $\mathcal{O}(10^{-90})$. This is a statement about the fidelity of the continuum limit, not about the truth of the discrete model.

\textbf{Flag:} The cell count $N\sim 10^{90}$ is a fiducial estimate. The actual number is not derived from A1-A3 and may differ by many orders of magnitude. However, as long as $N \gg 10^{20}$, the continuum approximation is effectively exact for all macroscopic purposes.

%%%%%%%%%%%%%%%%%%%%%%%%%%%%%%%%%%%%%%%%%%%%%%%%%%%%%%%%%%%%%%%%%%%%%
\section{Physical Precedents}
%%%%%%%%%%%%%%%%%%%%%%%%%%%%%%%%%%%%%%%%%%%%%%%%%%%%%%%%%%%%%%%%%%%%%

% ============================================================
\subsection{C1. Maxwell-Cattaneo Heat Conduction (Second Sound in NaF, He-4)}
% ============================================================

\subsubsection{C1.1 Historical context}

Fourier's law of heat conduction $\mathbf{J}_Q = -k\nabla T$ leads to the parabolic heat equation $\partial_t T = \alpha\nabla^2 T$, which predicts infinite propagation speed. Cattaneo (1948) and Vernotte (1958) independently proposed the modification:
\begin{equation}
\tau\,\partial_t\mathbf{J}_Q + \mathbf{J}_Q = -k\nabla T,
\end{equation}
which leads to the telegraph equation for temperature and predicts finite propagation speed $v = \sqrt{\alpha/\tau}$.

\subsubsection{C1.2 Experimental confirmation of second sound}

The wavelike propagation of heat (``second sound'') has been observed in:
\begin{itemize}
\item \textbf{Solid NaF} (Jackson \& Walker, 1971; Jackson, Walker, \& McNelly, 1970): Heat pulses propagated ballistically at low temperatures, with speed $\sim 2000$~m/s, confirming the Cattaneo prediction.
\item \textbf{Liquid He-4} (Peshkov, 1944): Second sound in superfluid helium at speeds $\sim 20$~m/s, though the microscopic mechanism (two-fluid model) differs from the Cattaneo relaxation.
\item \textbf{Bi and graphite} (Narayanamurti \& Dynes, 1972): Ballistic phonon propagation at low temperatures in dielectric crystals.
\end{itemize}

\subsubsection{C1.3 Structural analogy}

The Cattaneo heat equation and the Cattaneo $q$-field equation share the identical mathematical structure: a gradient-driven flux with finite relaxation time, producing a telegraph equation. The physics differs (thermal energy vs.\ information), but the mathematics is isomorphic. This provides a ``existence proof'' that the structure $\tau\partial_t\mathbf{J} + \mathbf{J} = -\kappa\nabla\phi$ is physically realizable in macroscopic systems.

\subsubsection{C1.4 Key difference}

In heat conduction, the Cattaneo correction matters only at low temperatures and high frequencies ($\omega\tau \gtrsim 1$). In the $q$-field framework, $\tau \sim \ell_P/c \sim 10^{-43}$~s, so the correction matters only at Planck frequencies $\omega \sim 10^{43}$~Hz. All laboratory and astrophysical phenomena are deep in the overdamped regime $\omega\tau \ll 1$, where the parabolic equation suffices.

\textbf{Reference:} D.D.\ Joseph \& L.\ Preziosi, Rev.\ Mod.\ Phys.\ \textbf{61}, 41 (1989); \textbf{62}, 375 (1990).

% ============================================================
\subsection{C2. London Equations (Superconductivity --- Current as Dynamical Variable)}
% ============================================================

\subsubsection{C2.1 The London equations}

The London brothers (1935) proposed that in a superconductor, the current density $\mathbf{J}_s$ is not merely a response to the electric field but a \textbf{dynamical variable} governed by:
\begin{equation}
\frac{\partial}{\partial t}(\Lambda\mathbf{J}_s) = \mathbf{E}, \qquad \Lambda = \frac{m}{n_s e^2},
\end{equation}
where $n_s$ is the superfluid density. Combined with the Maxwell equation $\nabla\times\mathbf{B} = \mu_0\mathbf{J}_s$, this yields:
\begin{equation}
\nabla^2\mathbf{B} = \frac{1}{\lambda_L^2}\mathbf{B},
\end{equation}
where $\lambda_L = \sqrt{m/(\mu_0 n_s e^2)}$ is the London penetration depth. The magnetic field is expelled from the superconductor (Meissner effect).

\subsubsection{C2.2 Current as a dynamical degree of freedom}

The crucial conceptual innovation of the London theory is elevating the current $\mathbf{J}_s$ from a constitutive relation ($\mathbf{J} = \sigma\mathbf{E}$) to a dynamical field with its own equation of motion. This is precisely the DGF move: $\mathbf{J}$ is elevated from $\mathbf{J} = -\kappa\nabla(\ln q)$ (constitutive) to a dynamical field obeying $\tau\partial_t\mathbf{J} + \mathbf{J} = -\kappa\nabla(\ln q)$.

\subsubsection{C2.3 Structural parallel}

\begin{center}
\begin{tabular}{lll}
\toprule
Concept & London Theory & DGF Cattaneo Extension \\
\midrule
Dynamical current & $\Lambda\partial_t\mathbf{J}_s = \mathbf{E}$ & $\tau\partial_t\mathbf{J} + \mathbf{J} = -\kappa\nabla(\ln q)$\\
Response time & $\tau_L = \Lambda/\rho_n$ (normal resistivity) & $\tau$ (Planck time)\\
Static limit & $\mathbf{J}_s$ screens $\mathbf{B}$ & $\mathbf{J}=0$ yields $\nabla^2(\ln q)=0$\\
Length scale & $\lambda_L$ (penetration depth) & $GM/c^2$ (gravitational radius)\\
\bottomrule
\end{tabular}
\end{center}

\textbf{Reference:} F.\ London \& H.\ London, Proc.\ R.\ Soc.\ A \textbf{149}, 71 (1935).

% ============================================================
\subsection{C3. Maxwell's Displacement Current ($\partial_t\mathbf{E}$ Term $\to$ Finite Propagation Speed)}
% ============================================================

\subsubsection{C3.1 Maxwell's correction to Amp\`ere's law}

Before Maxwell, Amp\`ere's law read $\nabla\times\mathbf{B} = \mu_0\mathbf{J}$. Taking the divergence gives $\nabla\cdot\mathbf{J}=0$, which contradicts the continuity equation $\nabla\cdot\mathbf{J} = -\partial_t\rho$ for time-varying charge distributions. Maxwell added the displacement current:
\begin{equation}
\nabla\times\mathbf{B} = \mu_0\mathbf{J} + \mu_0\epsilon_0\,\partial_t\mathbf{E}.
\end{equation}

The divergence now gives $\mu_0(\nabla\cdot\mathbf{J} + \epsilon_0\nabla\cdot\partial_t\mathbf{E}) = \mu_0(\nabla\cdot\mathbf{J} + \partial_t\rho) = 0$, consistent with charge conservation.

\subsubsection{C3.2 The $\partial_t$ term enables waves}

Maxwell's correction converts the electromagnetic equations from parabolic/elliptic to hyperbolic, enabling wave solutions with finite propagation speed $c = 1/\sqrt{\mu_0\epsilon_0}$. Without the $\partial_t\mathbf{E}$ term, electromagnetic disturbances would propagate instantaneously.

\subsubsection{C3.3 Precedent for the Cattaneo extension}

The DGF faces the identical mathematical situation. The parabolic $q$-field equation $\partial_t q = \kappa\nabla^2(\ln q)$ lacks a $\partial_t^2$ term and predicts infinite propagation speed. Adding the current relaxation $\tau\partial_t\mathbf{J}$ in the flux constitutive relation is analogous to Maxwell's addition of $\partial_t\mathbf{E}$:
\begin{itemize}
\item Maxwell: $0 = \nabla\times\mathbf{B} - \mu_0\mathbf{J}$ (static) $\to$ $\epsilon_0\partial_t\mathbf{E} = \nabla\times\mathbf{B} - \mu_0\mathbf{J}$ (dynamic).
\item DGF Cattaneo: $\mathbf{J} = -\kappa\nabla(\ln q)$ (static constitutive) $\to$ $\tau\partial_t\mathbf{J} + \mathbf{J} = -\kappa\nabla(\ln q)$ (dynamic constitutive).
\end{itemize}

In both cases, the $\partial_t$ term introduces a new fundamental constant ($\epsilon_0$ in EM, $\tau$ in DGF) and enables finite-speed propagation.

\textbf{Reference:} J.C.\ Maxwell, Phil.\ Trans.\ R.\ Soc.\ \textbf{155}, 459 (1865).

% ============================================================
\subsection{C4. Heaviside Transmission Line (Telegraph Equation Origin)}
% ============================================================

\subsubsection{C4.1 The telegrapher's equation}

Heaviside (1876) derived the equation for voltage $V(x,t)$ on a transmission line with distributed inductance $L$, capacitance $C$, resistance $R$, and leakage conductance $G$:
\begin{equation}
LC\,\partial_t^2 V + (LG+RC)\,\partial_t V + RG\,V = \partial_x^2 V.
\end{equation}

For a lossless line ($R=G=0$), this reduces to the wave equation $\partial_t^2 V = v^2\partial_x^2 V$ with $v=1/\sqrt{LC}$. For a line with $L=G=0$ (``Heaviside's distortionless'' limit), the equation is parabolic (pure diffusion). The general case with $L,R,C,G \neq 0$ is the telegrapher's equation---a damped wave equation.

\subsubsection{C4.2 The Cattaneo connection}

The Cattaneo heat equation and the telegrapher's equation share the same structure:
\begin{equation}
\partial_t^2 u + \gamma\partial_t u = v^2\nabla^2 u.
\end{equation}

Both can be derived from a flux with finite relaxation time. The telegrapher's equation is the ``hyperbolic regularization'' of the diffusion equation, and its derivation from a distributed-circuit model provides the physical intuition for the Cattaneo extension of the $q$-field: the cells in the DGF graph are analogous to the distributed $LC$ elements in a transmission line, with the finite propagation speed arising from the ``inductance'' (inertia) of the information current.

\subsubsection{C4.3 Mathematical properties}

The telegrapher's equation interpolates between the wave equation ($\gamma\to 0$) and the diffusion equation ($v^2/\gamma$ fixed, $\gamma\to\infty$). Its fundamental solution has compact support: disturbances propagate no faster than $v$. This is the key property that the Cattaneo extension confers on the $q$-field.

\textbf{Reference:} O.\ Heaviside, Phil.\ Mag.\ \textbf{2}, 135 (1876).

% ============================================================
\subsection{C5. Persistent Random Walks (Microscopic Foundation)}
% ============================================================

\subsubsection{C5.1 From random walk to diffusion}

A standard random walk in $d$ dimensions with step length $\Delta x$ and step time $\Delta t$ produces the diffusion equation $\partial_t p = D\nabla^2 p$ with $D = (\Delta x)^2/(2d\Delta t)$ in the continuum limit. This is the microscopic foundation of Fickian diffusion.

\subsubsection{C5.2 Persistent random walks}

A \textbf{persistent random walk} (F\"urth, 1917; Taylor, 1921; Goldstein, 1951) modifies the standard walk by introducing a correlation between successive steps: the walker has a probability $\alpha$ per step of continuing in the same direction and $1-\alpha$ of randomizing. Let $\mathbf{v}$ be the velocity and $p(\mathbf{x},\mathbf{v},t)$ the phase-space density. The governing kinetic equation is:
\begin{equation}
\partial_t p + \mathbf{v}\cdot\nabla p = -\frac{1}{\tau}p + \frac{1}{\tau}\rho(\mathbf{x},t)W(\mathbf{v}),
\end{equation}
where $\tau = \Delta t/(1-\alpha)$ is the velocity relaxation time, $\rho = \int p\,d\mathbf{v}$ is the spatial density, and $W(\mathbf{v})$ is the equilibrium velocity distribution.

Taking moments yields, for the spatial density, the \textbf{telegraph equation}:
\begin{equation}
\tau\partial_t^2\rho + \partial_t\rho = D\nabla^2\rho,
\end{equation}
with $D = \langle v^2\rangle\tau/d$. The propagation speed is finite: a persistent walker cannot move faster than its microscopic speed $v$.

\subsubsection{C5.3 Connection to DGF}

The DGF information current $\mathbf{J}$ corresponds to the net drift of ``information carriers'' (occupied cell states) on the graph. The finite relaxation time $\tau$ in the Cattaneo flux reflects the persistence of information movement: once information starts flowing in a direction, it tends to continue until scattered by capacity saturation. The persistent random walk provides the statistical-mechanical foundation for why the flux has inertia---it is the macroscopic manifestation of microscopic directional persistence.

\textbf{Reference:} G.H.\ Weiss, \emph{Aspects and Applications of the Random Walk} (North-Holland, 1994).

% ============================================================
\subsection{C6. Caldeira-Leggett Model (Wave $\to$ Diffusion Crossover from Bath Coupling)}
% ============================================================

\subsubsection{C6.1 The Caldeira-Leggett model}

Caldeira and Leggett (1983) studied a quantum particle coupled to a bath of harmonic oscillators. The reduced dynamics of the particle, after tracing out the bath, is described by a quantum Langevin equation. In the classical limit, the particle's probability density $P(x,t)$ evolves under the Fokker-Planck (Kramers) equation:
\begin{equation}
\partial_t P + v\partial_x P - \frac{1}{m}V'(x)\partial_v P = \gamma\partial_v(vP) + \frac{\gamma k_B T}{m}\partial_v^2 P.
\end{equation}

In the high-friction (Smoluchowski) limit $\gamma \to \infty$, this reduces to the Smoluchowski (diffusion) equation:
\begin{equation}
\partial_t\rho = \frac{1}{m\gamma}\partial_x(V'(x)\rho) + \frac{k_B T}{m\gamma}\partial_x^2\rho,
\end{equation}
which is parabolic. At finite $\gamma$, the dynamics is hyperbolic (Kramers equation has a $\partial_t^2$ term after eliminating $v$). The crossover from wavelike to diffusive behavior is controlled by the friction coefficient $\gamma$.

\subsubsection{C6.2 Structure of the crossover}

In the underdamped regime ($\gamma$ small), the particle oscillates in the potential well with slowly decaying amplitude---wavelike behavior. In the overdamped regime ($\gamma$ large), the particle drifts toward the potential minimum without oscillation---diffusive behavior. The crossover time is $\tau_{\rm cross} \sim 1/\gamma$.

\subsubsection{C6.3 Precedent for the DGF wave-diffusion crossover}

The DGF Cattaneo system has the identical mathematical structure: the damping rate $\gamma = 1/\tau$ controls the crossover from wave-like ($\omega\tau \gg 1$) to diffusive ($\omega\tau \ll 1$) behavior. The Caldeira-Leggett model provides the microscopic justification: a system coupled to a bath (environment) generically exhibits wave-to-diffusion crossover, with the diffusion limit emerging when the bath coupling is strong. In the DGF context, the ``bath'' is the ensemble of all other cells, and the coupling is the information flux.

\textbf{Reference:} A.O.\ Caldeira \& A.J.\ Leggett, Ann.\ Phys.\ \textbf{149}, 374 (1983).

%%%%%%%%%%%%%%%%%%%%%%%%%%%%%%%%%%%%%%%%%%%%%%%%%%%%%%%%%%%%%%%%%%%%%
\section{Gravitational Wave Details}
%%%%%%%%%%%%%%%%%%%%%%%%%%%%%%%%%%%%%%%%%%%%%%%%%%%%%%%%%%%%%%%%%%%%%

\textbf{Preamble:} The current $q$-field equation $\partial_t q = \kappa\nabla^2(\ln q)$ is parabolic and cannot describe propagating gravitational waves. The following analysis assumes the hyperbolic extension $\partial_t^2 q + \gamma\partial_t q = c_s^2\nabla^2(\ln q)$ (with $\gamma = 1/\tau$, $c_s^2 = \kappa/\tau$), which reduces to the parabolic equation in the static/overdamped limit. All predictions in this section depend on the hyperbolic extension and its coupling to matter, neither of which is rigorously derived from A1-A3. The material below represents a research program, not established results.

% ============================================================
\subsection{D1. Derivation of Source Term $S(\mathbf{x},t)$ from Binary $q$-Field Quadrupole}
% ============================================================

\subsubsection{D1.1 Scalar wave equation with source}

Generalizing the Poisson equation $\nabla^2\psi = -(4\pi G/c^2)\rho_m$ to the time-dependent case, the simplest Lorentz-invariant extension is:
\begin{equation}
\Box\psi \equiv \left(-\frac{1}{c^2}\partial_t^2 + \nabla^2\right)\psi = -\frac{4\pi G}{c^2}\,\rho_m(\mathbf{x},t),
\label{eq:scalar_wave}
\end{equation}
where $\Box$ is the d'Alembertian operator and $\psi = -\ln q$.

\subsubsection{D1.2 Green's function and retarded solution}

The retarded Green's function for the wave equation in 3+1 dimensions is:
\begin{equation}
G_{\rm ret}(\mathbf{x},t;\mathbf{x}',t') = \frac{\delta(t - t' - |\mathbf{x}-\mathbf{x}'|/c)}{4\pi|\mathbf{x}-\mathbf{x}'|}.
\end{equation}

The solution is:
\begin{equation}
\psi(\mathbf{x},t) = \frac{G}{c^2}\int \frac{\rho_m(\mathbf{x}', t - |\mathbf{x}-\mathbf{x}'|/c)}{|\mathbf{x}-\mathbf{x}'|}\,d^3x'.
\end{equation}

Far from the source ($r \gg$ source size), expand $|\mathbf{x}-\mathbf{x}'| \approx r - \mathbf{x}'\cdot\hat{\mathbf{n}}$ where $r=|\mathbf{x}|$, $\hat{\mathbf{n}} = \mathbf{x}/r$:
\begin{equation}
\psi(\mathbf{x},t) \approx \frac{G}{c^2 r}\int \rho_m(\mathbf{x}', t - r/c + \mathbf{x}'\cdot\hat{\mathbf{n}}/c)\,d^3x'.
\end{equation}

\subsubsection{D1.3 Multipole expansion}

For a non-relativistic source ($v \ll c$), expand the retarded time:
\begin{equation}
\rho_m(\mathbf{x}', t - r/c + \mathbf{x}'\cdot\hat{\mathbf{n}}/c) \approx \rho_m(\mathbf{x}', t_{\rm ret}) + \frac{\mathbf{x}'\cdot\hat{\mathbf{n}}}{c}\,\partial_t\rho_m(\mathbf{x}', t_{\rm ret}) + \cdots,
\end{equation}
where $t_{\rm ret} = t - r/c$.

\textbf{Monopole:} $\int\rho_m\,d^3x' = M$ (constant). This gives the static term $\psi \sim GM/c^2 r$, which does not radiate (no time dependence).

\textbf{Dipole:} $\int \mathbf{x}'\rho_m\,d^3x' = M\mathbf{X}_{\rm CM}(t_{\rm ret})$. For an isolated system, $\mathbf{X}_{\rm CM} = \text{const}$ (conservation of momentum), so the dipole term vanishes identically.

\textbf{Quadrupole:} The leading radiating term is the mass quadrupole:
\begin{equation}
\boxed{\psi(\mathbf{x},t) \approx \frac{G}{2c^4 r}\,\hat{n}^i\hat{n}^j\,\ddot{I}_{ij}(t_{\rm ret})},
\label{eq:quadrupole}
\end{equation}
where $I_{ij}(t) \equiv \int \rho_m(\mathbf{x},t)\,x^i x^j\,d^3x$ is the mass quadrupole moment tensor. The factor $1/c^4$ arises from two time derivatives (each $1/c^2$) applied to $I_{ij}$.

\subsubsection{D1.4 Binary source}

For a binary system of masses $m_1, m_2$ in circular orbit with separation $R$ and orbital frequency $\Omega = \sqrt{G(m_1+m_2)/R^3}$, the quadrupole moment oscillates at frequency $2\Omega$:
\begin{equation}
\ddot{I}_{ij} \propto \mu R^2\Omega^2\,e^{2i\Omega t_{\rm ret}},
\end{equation}
where $\mu = m_1 m_2/(m_1+m_2)$ is the reduced mass. The $\psi$-field radiation has frequency $f = \Omega/\pi$.

\textbf{Key difference from GR:} In GR, the gravitational wave strain $h$ from a binary has both $+$ and $\times$ tensor polarizations, and the amplitude scales as $(G\mathcal{M}/c^3)^{5/3}/r$ where $\mathcal{M}$ is the chirp mass. In scalar DGF, the radiation is in the scalar field $\psi = -\ln q$, with a single (breathing) polarization, and the amplitude scales as $G\mu R^2\Omega^2/(c^4 r)$.

\textbf{Approximation flagged:} This entire derivation assumes (i) the scalar wave equation (\ref{eq:scalar_wave}) as the relativistic extension, (ii) minimal coupling to matter via $\rho_m$, and (iii) the quadrupole formula for non-relativistic sources. None of these are derived from A1-A3. They represent the simplest Lorentz-invariant generalization of the static Poisson equation.

% ============================================================
\subsection{D2. Scalar Polarization Pattern (Breathing Mode B) vs.\ GR Tensor Modes ($+$, $\times$)}
% ============================================================

\subsubsection{D2.1 GR tensor polarizations}

In general relativity, gravitational waves are transverse-traceless tensor perturbations $h_{ij}^{\rm TT}$ of the metric, with two independent polarization states:
\begin{align}
h_+ &= A_+ \cos(2\phi)\cos(\omega t - kz),\\
h_\times &= A_\times \sin(2\phi)\cos(\omega t - kz),
\end{align}
where $\phi$ is the polarization angle. Under a rotation by $\pi/2$, $h_+ \to -h_+$ and $h_\times \to -h_\times$ (spin-2). The effect on a ring of test particles is a quadrupolar deformation: $+$ mode stretches along $x$ and compresses along $y$; $\times$ mode does the same along the $45^\circ$ axes.

\subsubsection{D2.2 Scalar (breathing) polarization}

A scalar gravitational wave (as predicted by scalar-tensor theories and by the DGF $\psi$-field) produces isotropic strain: a ring of test particles expands and contracts uniformly. The polarization pattern is:
\begin{equation}
h_B(t) = A_B \cos(\omega t - kz),
\end{equation}
independent of $\phi$. Under rotations, $h_B$ is invariant (spin-0).

\subsubsection{D2.3 Detector response}

A gravitational wave detector with arm vectors $\hat{\mathbf{u}}$, $\hat{\mathbf{v}}$ measures the strain:
\begin{equation}
h(t) = \frac{1}{2}(u^i u^j - v^i v^j)h_{ij}^{\rm TT}.
\end{equation}

For the GR $+$,$\times$ modes: the antenna pattern functions are:
\begin{align}
F_+(\theta,\phi,\psi) &= \frac{1}{2}(1+\cos^2\theta)\cos 2\phi\cos 2\psi - \cos\theta\sin 2\phi\sin 2\psi,\\
F_\times(\theta,\phi,\psi) &= \frac{1}{2}(1+\cos^2\theta)\cos 2\phi\sin 2\psi + \cos\theta\sin 2\phi\cos 2\psi.
\end{align}

For the scalar breathing mode, the antenna pattern is:
\begin{equation}
F_B(\theta,\phi) = -\frac{1}{2}\sin^2\theta\cos 2\phi,
\end{equation}
which has a distinctly different angular dependence from the tensor modes.

\subsubsection{D2.4 Multi-detector discrimination}

With three or more non-co-located detectors (e.g., the LIGO-Virgo network), the different antenna patterns allow model selection between pure tensor (GR), pure scalar (DGF-like), and mixed polarization hypotheses. The null stream method (G\"ursel \& Tinto, 1989) can isolate scalar contributions even in the presence of dominant tensor signals.

\textbf{Reference:} C.M.\ Will, Living Rev.\ Relativ.\ \textbf{17}, 4 (2014); B.P.\ Abbott et al.\ (LIGO/Virgo), Phys.\ Rev.\ Lett.\ \textbf{116}, 061102 (2016).

% ============================================================
\subsection{D3. LIGO/Virgo Polarization Analysis Methods}
% ============================================================

\subsubsection{D3.1 Bayesian model selection}

LIGO/Virgo test polarization content by comparing the Bayesian evidence for different polarization models:
\begin{equation}
\mathcal{B}_{S/T} = \frac{P(d|\text{tensor})}{P(d|\text{scalar})}\cdot\frac{P(\text{tensor})}{P(\text{scalar})}.
\end{equation}

For GW170817, the log Bayes factor favors pure tensor over pure scalar by $\ln\mathcal{B} \sim 20$ (overwhelming evidence for tensor). For pure tensor over pure vector: $\ln\mathcal{B} \sim 18$. These constraints assume a single polarization family; mixed models (tensor + scalar admixture) are less well constrained.

\subsubsection{D3.2 Null stream analysis}

With $N$ detectors, there are $N-2$ independent null streams for tensor waves and $N-1$ for scalar waves. The null stream energy is:
\begin{equation}
E_{\rm null} = \sum_{I=1}^{N_{\rm null}} \int |\tilde{n}_I(f)|^2\,df,
\end{equation}
where $\tilde{n}_I(f)$ is the Fourier transform of the $I$-th null stream. A statistically significant non-zero $E_{\rm null}$ for the tensor null streams, combined with zero $E_{\rm null}$ for the scalar null stream, would indicate a scalar polarization component.

\subsubsection{D3.3 Current constraints from GW170817}

Joint LIGO-Virgo analysis of GW170817 (Phys.\ Rev.\ Lett.\ \textbf{123}, 011102, 2019):
\begin{itemize}
\item Pure tensor: favored over pure scalar by $\ln\mathcal{B} > 20$.
\item Pure tensor: favored over pure vector by $\ln\mathcal{B} > 18$.
\item Allowed scalar amplitude: $< 25\%$ of tensor amplitude (90\% credible upper limit).
\end{itemize}

A DGF scalar wave with amplitude comparable to the GR tensor amplitude would have been detected. If the scalar-to-tensor amplitude ratio is below $\sim 25\%$, it remains consistent with current data.

\textbf{Reference:} LIGO/Virgo Collaboration, Phys.\ Rev.\ Lett.\ \textbf{123}, 011102 (2019).

% ============================================================
\subsection{D4. Predicted Waveform for Scalar GW from Binary Merger}
% ============================================================

\subsubsection{D4.1 Inspiral waveform}

In the stationary phase approximation, the Fourier-domain scalar waveform for a binary inspiral is:
\begin{equation}
\tilde{h}_B(f) = A_B\,f^{-7/6}\,e^{i\Psi_B(f)},
\end{equation}
where the amplitude $A_B \propto \mathcal{M}_c^{5/3}/D_L$ (same mass and distance scaling as GR), and the phase:
\begin{equation}
\Psi_B(f) = 2\pi f t_c - \phi_c - \frac{\pi}{4} + \frac{3}{128}(\pi\mathcal{M}_c f)^{-5/3}\left[1 + \Delta\Psi_{\rm PN}(f)\right],
\end{equation}
with post-Newtonian corrections $\Delta\Psi_{\rm PN}$.

\subsubsection{D4.2 Key differences from GR}

\begin{enumerate}[label=(\roman*)]
\item \textbf{Frequency evolution:} Scalar radiation carries energy away from the binary, causing the orbit to decay. The rate of energy loss to scalar radiation differs from GR's quadrupole formula. This shifts the chirp rate $\dot{f}$:
   \begin{equation}
   \dot{f}_{\rm scalar} = \dot{f}_{\rm GR}\times(1 + \delta),
   \end{equation}
   where $\delta$ depends on the scalar-tensor coupling strength. For Brans-Dicke theory with $\omega_{\rm BD}=0$ (maximal scalar coupling), $\delta \sim +1/3$ (faster chirp). For $\omega_{\rm BD} \gg 1$, $\delta \to 0$.

\item \textbf{Phase evolution:} The accumulated phase difference over the LIGO band ($\sim 10$--$1000$~Hz) depends on the number of cycles:
   \begin{equation}
   \Delta\Phi = 2\pi\int_{f_{\rm low}}^{f_{\rm high}} \frac{f}{\dot{f}}\,df.
   \end{equation}
   A $1\%$ shift in $\dot{f}$ over $\sim 100$ cycles produces $\sim 1$~radian phase shift---detectable with matched filtering.

\item \textbf{Polarization content:} Pure breathing mode vs.\ $+,\times$ tensor modes---detectable via antenna pattern analysis.
\end{enumerate}

\subsubsection{D4.3 Merger and ringdown}

The DGF framework currently has no description of the merger phase (strong field, relativistic velocities) or the ringdown (requires a wave equation for the compact object). In GR, the post-merger signal is a superposition of quasi-normal modes; in DGF, the absence of a sharp event horizon would modify the QNM spectrum, but no quantitative prediction exists.

\textbf{Flag:} The inspiral waveform above is derived from the scalar wave equation (\ref{eq:scalar_wave}) and quadrupole formula (\ref{eq:quadrupole}). It assumes the $\psi$-field couples to matter with strength $G/c^2$ (derived from the static Newtonian calibration). The actual DGF waveform depends on the hyperbolic extension and the coupling mechanism, both of which are open problems. The inspiral waveform is presented as a target for theory development, not a prediction.

% ============================================================
\subsection{D5. GW170817 Constraint Analysis (Speed = $c$ to $10^{-15}$)}
% ============================================================

\subsubsection{D5.1 The multimessenger event}

GW170817 (August 17, 2017) was a binary neutron star merger detected in gravitational waves by LIGO-Virgo and in gamma rays by Fermi-GBM and INTEGRAL. The GRB 170817A was observed $1.74 \pm 0.05$~s after the GW merger time.

\subsubsection{D5.2 Speed of gravitational waves}

The arrival time difference constrains the difference between the speed of gravitational waves $v_{\rm GW}$ and the speed of light $c$:
\begin{equation}
\frac{|v_{\rm GW} - c|}{c} \le 10^{-15}.
\end{equation}

This is a direct consequence of the $1.74$~s delay over a distance of $\sim 40$~Mpc ($\sim 1.2\times 10^{24}$~m). The light travel time is $\sim 1.3\times 10^8$ years, so a fractional speed difference of $10^{-15}$ would produce a relative delay of:
\begin{equation}
\Delta t = 1.3\times 10^8~\text{yr} \times 10^{-15} \times 3.15\times 10^7~\text{s/yr} \sim 4~\text{s},
\end{equation}
which would have been detected.

\subsubsection{D5.3 Implication for the $q$-field hyperbolic extension}

If DGF gravitational waves are described by the hyperbolic extension (\ref{eq:telegraph}), their propagation speed is:
\begin{equation}
c_s = \sqrt{\frac{\kappa}{\tau q_0}}.
\end{equation}

GW170817 requires $c_s = c$ to within $10^{-15}$. This constrains the microscopic parameters:
\begin{equation}
\frac{\kappa}{\tau} = c^2 q_0 \quad \text{to within } 10^{-15}.
\end{equation}

Since $q_0 \approx 1$ in the intergalactic medium (nearly pure vacuum), this is effectively a constraint on the ratio $\kappa/\tau$. The DGF framework does not predict this ratio independently; it must be calibrated to match the observed $c$.

\subsubsection{D5.4 What this means for DGF}

\begin{itemize}
\item The hyperbolic extension \textbf{must} have $c_s = c$ to satisfy GW170817. This is not an automatic consequence of A1-A3---it is an additional calibration.
\item The parabolic equation also has infinite propagation speed, which technically violates GW170817. However, the ``infinite speed'' is a continuum artifact: the underlying discrete dynamics operates at finite speed $c$ (by A3). The observed finite $c$ for gravitational waves is a consistency check that the DGF communication bound is the same as the gravitational wave propagation speed.
\item If $c_s \neq c$, scalar DGF waves would arrive before/after the gamma rays, producing a measurable time offset. The $10^{-15}$ constraint effectively rules out any DGF model where $c_s$ differs from $c$ by more than one part in $10^{15}$.
\end{itemize}

\textbf{Reference:} LIGO/Virgo Collaboration, Astrophys.\ J.\ \textbf{848}, L13 (2017).

%%%%%%%%%%%%%%%%%%%%%%%%%%%%%%%%%%%%%%%%%%%%%%%%%%%%%%%%%%%%%%%%%%%%%
\section{Inspector Audit Trail}
%%%%%%%%%%%%%%%%%%%%%%%%%%%%%%%%%%%%%%%%%%%%%%%%%%%%%%%%%%%%%%%%%%%%%

% ============================================================
\subsection{E1. List of All Issues Found by Internal Audit}
% ============================================================

The DGF framework underwent nine rounds of adversarial review (REVIEWER A through I) plus multiple INSPECTOR verification rounds. The following is the complete issue inventory.

\begin{table}[h]
\caption{Complete audit issue inventory. Severity: FATAL (must fix), SEVERE (should fix), MODERATE (consider), MINOR (cosmetic).}
\begin{tabular}{p{0.5cm}p{7cm}l}
\toprule
ID & Issue & Severity \\
\midrule
G1 & Independence circularity: Shannon additivity requires independent cells; flux creates correlations. Derivation not self-consistent. & FATAL \\
G6 & Uniqueness-vs-universality trilemma: cannot claim both that $\Phi=-\ln q$ is uniquely selected AND that all $\varphi\in\Phi$ give same physics AND that specific predictions distinguish the theory. & FATAL \\
H2 & Shadow prediction derived from ad hoc refractive index model, not from DGF. & FATAL \\
H6 & DGF shadow prediction excluded by EHT M87* at 7$\sigma$ in the only quantitative model. & FATAL \\
H1 & Parabolic equation cannot describe gravitational waves; ringdown predictions unfounded. & FATAL (self-acknowledged) \\
G5b & Shannon citation inflation: ``selected by Shannon's uniqueness theorem'' overstates. & SEVERE \\
G2 & Bridge from Shannon entropy to flux potential is a physical postulate (B1), not derived. & SEVERE \\
G4 & Irrotationality ($\oint J=0$) and gradient flow assumptions not in assumption inventory. & SEVERE \\
G5a & B1 (flux potential must be information measure) not acknowledged as assumption. & SEVERE \\
H4 & One quantitative ringdown estimate (3$\times$ frequency shift) self-admittedly excluded by LIGO. & SEVERE \\
H5 & Information paradox ``resolution'': replacing ``causally impossible'' with $e^{-10^{76}}$ is relabeling. & SEVERE \\
A8 & Entropy identity: entropy can decrease between steps (not a Clausius monotone). & SEVERE \\
I1 & Systematic conflation of ``consistent with'' and ``forced by'' throughout the paper. & SEVERE \\
H3 & Physical state of matter at $q\to0$ cores undefined---what happens to infalling matter? & SEVERE \\
C4 & Sandpile (BTW) analogy for self-organized criticality not quantitatively accurate. & MODERATE \\
G3 & Choice of vacant ($-\ln q$) vs.\ occupied ($-\ln(1-q)$) self-information: physically defensible but argument unclear. & MODERATE \\
G7 & C5 (vacuum stability) is a substantive physical postulate, not a definition. & MODERATE \\
B8 & No experimental prediction that is simultaneously (a) quantitative, (b) derived from DGF, and (c) not already excluded. & MODERATE \\
D8 & Title ``Why the Universe Looks Different at Different Scales'' contradicts Honest Boundaries section. & MODERATE \\
E1 & Flux function chosen, not derived---acknowledged in v3 but narrative still implies necessity. & MODERATE \\
C6 & Cosomological inflation explanation: scale factor $N_e \sim \ln(q_f/q_i)$ requires plausibility argument for $q_i$. & MINOR \\
\bottomrule
\end{tabular}
\end{table}

% ============================================================
\subsection{E2. Resolution of Each Issue}
% ============================================================

\begin{enumerate}[leftmargin=*, label=\textbf{G1:}]
\item \textbf{Independence circularity.} \textit{Resolution:} The derivation is reframed as a modeling ansatz: $\Phi=-\ln q$ is the flux potential in the limit of uncorrelated cells (the ``ideal flux''). This is analogous to the ideal gas law---correct in the non-interacting limit, with corrections when interactions (flux) are present. The universality theorem guarantees that qualitative conclusions are insensitive to correlation corrections. This is acknowledged explicitly in the main text as an approximation, not a deduction from Shannon.

\textit{Remaining concern:} The formal derivation of correlation corrections to the ideal flux is an open problem.

\item \textbf{G6: Uniqueness-vs-universality trilemma.} \textit{Resolution:} The paper now cleanly separates (a) qualitative conclusions that are universal across $\Phi$ (dual-domain, Laplace static limit, $1/r$ decay, no sharp horizon), from (b) quantitative predictions that depend on the specific $\varphi$ choice. The table in Attack 6 analysis (Reviewer G) is incorporated. Shannon's theorem is presented as \textbf{motivation} for the specific choice $\Phi=-\ln q$, not as a \textbf{derivation} of uniqueness.

\item \textbf{H2 + H6: Decorative shadow prediction.} \textit{Resolution:} The black hole shadow prediction ($\sim 2.72\,GM/c^2$) is removed from the main paper. It was computed in an ad hoc refractive index model ($n=1/q$) that is not derived from DGF principles and is excluded by EHT M87* at 7$\sigma$. The main paper now states that quantitative EHT predictions require (i) an effective metric $g_{\mu\nu}[q]$ and (ii) a field equation for light propagation in the $q$-field background---neither of which exists.

\item \textbf{H1 + H4: Ringdown prediction unfounded.} \textit{Resolution:} The claim that ``DGF ringdown differs from GR'' is removed. The main paper acknowledges that the parabolic equation cannot describe gravitational waves and that the hyperbolic extension is required. Ringdown predictions are demoted from ``observational signatures'' to ``future research targets.''

\item \textbf{G2 + G5a: B1 not acknowledged.} \textit{Resolution:} The assumption inventory is expanded to include B1 (``$\Phi$ must be an additive information measure'') as an explicit postulate. The derivation chain from Shannon's theorem to $J=-\kappa\nabla(\ln q)$ is presented with all required postulates enumerated.

\item \textbf{G4 + G7: Missing postulates in assumption inventory.} \textit{Resolution:} The assumption inventory now lists all postulates: A1, A2, A3, R1 (irrotationality), R2 (gradient flow), B1 (information measure), B2 ($q$ as probability), C5 (vacuum stability), $\kappa$, $\eta$, $d=3$. Total: 11 assumptions (3 definitions + 8 inputs).

\item \textbf{H5: Information paradox ``resolution.''} \textit{Resolution:} The claim of solving the information paradox is downgraded. The main paper now states: ``DGF preserves unitarity in principle because $f$ is a permutation on a finite set. The probability of information reflux for a macroscopic black hole is $P_{\rm reflux} \le e^{-N_S} \sim e^{-10^{76}}$. Whether a probability this small constitutes a resolution of the information paradox is a matter of interpretation on which reasonable physicists differ.''

\item \textbf{I1: Conflation of ``consistent with'' and ``forced by.''} \textit{Resolution:} A systematic language audit was performed. ``We prove that...'' is replaced by ``We show that...'' for conclusions that depend on modeling choices. ``Derived from A1-A3'' is replaced by ``constructed from A1-A3 together with the following postulates...'' The distinction between ``the axioms force this conclusion'' and ``the axioms are consistent with this conclusion'' is maintained throughout.

\item \textbf{H3: Physical state at $q\to0$.} \textit{Resolution:} Acknowledged as an open problem. The discrete DGF model defines the state as a bit-string on $N$ cells; $q\to0$ means almost all cells in the core region are occupied ($s_i=1$). The continuum limit breaks down at the cell scale. The question ``what is the matter density at $r\to0$'' is ill-posed in the continuum approximation and requires the discrete formulation.

\item \textbf{A8: Entropy can decrease.} \textit{Resolution:} The entropy identity in the main paper is presented with the explicit caveat: entropy can decrease between consecutive time steps while remaining above its initial value. This is not a violation of the second law---it is a feature of exact microscopic dynamics (fluctuation theorem structure). On macroscopic scales ($N$ large), decreases are exponentially rare.
\end{enumerate}

% ============================================================
\subsection{E3. Remaining Open Problems with Honest Assessment}
% ============================================================

\begin{enumerate}[leftmargin=*, label=\textbf{OP-\arabic*:}]

\item \textbf{Cell basis problem.} What selects the $\{|0\rangle,|1\rangle\}$ basis in which cells are ``occupied'' or ``vacant''? Conjecture: the basis is selected by interaction history---cells that repeatedly participate in overflow events have their basis locked. \textbf{Status: Unproved conjecture.} This is the deepest open problem. Until resolved, the $q$-field is defined relative to an externally specified preferred basis.

\item \textbf{Hyperbolic extension.} The $q$-field needs a hyperbolic (wave-type) extension to describe gravitational waves and finite-speed propagation. Candidate: $\partial_t^2 q + \gamma\partial_t q = c_s^2\nabla^2(\ln q)$. \textbf{Status: Postulated, not derived.} No derivation of this specific form from A1-A3 exists. The parameters $\gamma$ and $c_s$ are additional free parameters.

\item \textbf{Coupling to matter and radiation.} How does the $q$-field couple to matter and electromagnetic fields? Is there an effective metric $g_{\mu\nu}[q]$? \textbf{Status: Completely open.} Five possible coupling schemes have been enumerated (refractive index, conformal metric, Brans-Dicke, Finsler geometry, no coupling); they give different predictions. No first-principles derivation exists.

\item \textbf{Strong-field regime.} The exact solution to $\nabla^2(\ln q)=0$ with boundary condition $q(r_s)\to0$ is not known in closed form. The weak-field expansion $q\approx 1-GM/rc^2$ breaks down at $r\sim GM/c^2$. \textbf{Status: Conjectured but not proved.} The full solution may require a discrete treatment.

\item \textbf{Spatial dimensionality.} DGF does not predict $d=3$. The $1/r$ decay of gravity and the dual-domain structure both assume $d=3$. In $d\neq 3$, the predictions change qualitatively. \textbf{Status: $d=3$ is an observational input.}

\item \textbf{Numerical constants: $G$, $c$, $\hbar$.} The coupling $\kappa = 4\pi G/c^2$ is calibrated by matching to Newtonian gravity. Neither $G$, $c$, nor $\hbar$ is derived from A1-A3. \textbf{Status: Acknowledged calibration.} The framework explains the \emph{form} of the gravitational field equation but not the \emph{value} of the coupling constant.

\item \textbf{Born rule.} DGF does not derive the Born rule (quantum measurement probabilities). Standard quantum mechanics requires it as an additional postulate. \textbf{Status: Not addressed.} The $q$-field describes the capacity for quantum coherence; it does not determine which specific outcome occurs in a measurement.

\item \textbf{Correlation corrections to the flux.} The self-information flux $\Phi=-\ln q$ assumes uncorrelated cells (Shannon additivity). Flux creates correlations, modifying the effective potential. \textbf{Status: Open.} The correction $\Phi_{\rm corr}(q_i,q_j) = -\ln q_i - \ln q_j - I(q_i:q_j)$ involves the mutual information, which depends on the flux history. This is a self-consistency problem.

\item \textbf{Discrete-to-continuum limit for the hyperbolic extension.} The discrete dynamics $f:X\to X$ operates in discrete time at finite speed $c$. Deriving the hyperbolic PDE from the discrete permutation dynamics---including the specific values of $\tau$ and $c_s$---is an open combinatorial problem. \textbf{Status: Not addressed.}

\item \textbf{Cosmology.} DGF has no cosmological extension. The $q$-field equation on a flat background does not incorporate expansion, and the framework does not address dark energy, the cosmological constant, or structure formation beyond qualitative arguments. \textbf{Status: Open research direction.}

\end{enumerate}

\textbf{Honest assessment:} Of the 10 open problems, OP-1 (cell basis) and OP-2 (hyperbolic extension) are the most urgent for the framework's coherence. OP-3 (coupling to matter) is essential for any empirical contact beyond the static Newtonian correspondence. The remaining problems are important but less immediately blocking.

%%%%%%%%%%%%%%%%%%%%%%%%%%%%%%%%%%%%%%%%%%%%%%%%%%%%%%%%%%%%%%%%%%%%%
\section{Numerical Estimates}
%%%%%%%%%%%%%%%%%%%%%%%%%%%%%%%%%%%%%%%%%%%%%%%%%%%%%%%%%%%%%%%%%%%%%

% ============================================================
\subsection{F1. $\gamma$ Estimate from Cosmological Constraints}
% ============================================================

\subsubsection{F1.1 Definition}

The damping rate $\gamma = 1/\tau$ in the Cattaneo extension. We estimate $\gamma$ from two constraints:
\begin{enumerate}[label=(\roman*)]
\item \textbf{Propagation speed:} $c_s^2 = \kappa\gamma/q_0$ must equal $c^2$ (from GW170817).
\item \textbf{Static calibration:} $\kappa$ is calibrated to reproduce Newtonian gravity in the static limit.
\end{enumerate}

\subsubsection{F1.2 Planck-scale estimate}

If the fundamental cell spacing is the Planck length $a = \ell_P = \sqrt{\hbar G/c^3} \approx 1.6\times 10^{-35}$~m, and the information propagation speed is $c$, then the natural timescale for a single cell update is:
\begin{equation}
\tau_0 = \frac{a}{c} = \frac{\ell_P}{c} \approx 5.4\times 10^{-44}~\text{s}.
\end{equation}

The diffusivity in the discrete model is $\kappa \sim a^2/\tau_0 = a c = \ell_P c \approx 4.8\times 10^{-27}$~m$^2$/s.

From the calibration $\kappa = 4\pi G/c^2$ (to match Newtonian gravity):
\begin{equation}
\frac{4\pi G}{c^2} = \frac{4\pi \times 6.67\times 10^{-11}}{9\times 10^{16}} \approx 9.3\times 10^{-27}~\text{m}^2/\text{s}.
\end{equation}

Comparing with $\kappa \sim \ell_P c$:
\begin{equation}
\ell_P c = \sqrt{\frac{\hbar G}{c^3}}\cdot c = \sqrt{\frac{\hbar G}{c}} \approx 4.8\times 10^{-27}~\text{m}^2/\text{s}.
\end{equation}

The factor of $\sim 2$ discrepancy ($9.3$ vs.\ $4.8$) is of order unity and can be absorbed into the $\mathcal{O}(1)$ prefactor in the cell dynamics. The Planck-scale identification is dimensionally consistent.

\subsubsection{F1.3 $\gamma$ from $c_s = c$}

Taking $c_s^2 = \kappa\gamma/q_0 = c^2$ with $q_0\approx 1$ (near-vacuum):
\begin{equation}
\gamma = \frac{c^2}{\kappa} \approx \frac{9\times 10^{16}}{9.3\times 10^{-27}} \approx 10^{43}~\text{s}^{-1}.
\end{equation}

The corresponding relaxation time:
\begin{equation}
\tau = \frac{1}{\gamma} \approx 10^{-43}~\text{s},
\end{equation}
which is the Planck time. This is consistent: the natural timescale for the information current to relax is the Planck time.

\subsubsection{F1.4 Alternative estimate from nuclear scale}

If the cell spacing is the nuclear scale $a \sim 10^{-15}$~m (as suggested by the $N\sim 10^{90}$ estimate), then $\tau_0 = a/c \sim 3\times 10^{-24}$~s, $\kappa \sim a c \sim 3\times 10^{-7}$~m$^2$/s, and $\gamma \sim 3\times 10^{23}$~s$^{-1}$. The calibration $\kappa = 4\pi G/c^2 \sim 10^{-26}$~m$^2$/s is inconsistent with $a \sim 10^{-15}$~m by a factor of $\sim 10^{19}$.

\textbf{Conclusion:} The Planck-scale identification is required for dimensional consistency between $\kappa$, $c$, and $G$. The cell spacing must be of order the Planck length (or a factor $\sim 2$ thereof). The ``nuclear scale'' $N\sim 10^{90}$ estimate in \S B5 is wrong for the cell scale (it was an estimate for the number of quantum field theory degrees of freedom, not the number of DGF cells). The correct DGF cell count is $N \sim (c/H_0\ell_P)^3 \sim 10^{184}$.

\textbf{Flag:} This dimensional consistency check is suggestive but does not constitute a derivation of $\ell_P$. The identification $a\sim\ell_P$ is an input, not an output.

% ============================================================
\subsection{F2. Diffusion-to-Wave Crossover Wavenumber $k_{\rm cross} = \gamma/(2c)$}
% ============================================================

\subsubsection{F2.1 Derivation}

From the dispersion relation for the Cattaneo extension (\S B2.3):
\begin{equation}
\sigma_\pm = -\frac{\gamma}{2} \pm \sqrt{\frac{\gamma^2}{4} - c^2 k^2},
\end{equation}
where $k = |\mathbf{k}|$, $c = c_s$, and $\gamma = 1/\tau$.

The crossover from real (overdamped, diffusive) to complex (underdamped, wavelike) eigenvalues occurs when the discriminant vanishes:
\begin{equation}
\frac{\gamma^2}{4} - c^2 k_{\rm cross}^2 = 0 \quad\Longrightarrow\quad \boxed{k_{\rm cross} = \frac{\gamma}{2c} = \frac{1}{2c\tau}}.
\label{eq:kcross}
\end{equation}

\subsubsection{F2.2 Numerical value}

With $\gamma \sim 10^{43}$~s$^{-1}$ (Planck-scale estimate) and $c = 3\times 10^8$~m/s:
\begin{equation}
k_{\rm cross} \approx \frac{10^{43}}{2\times 3\times 10^8} \approx 1.7\times 10^{34}~\text{m}^{-1}.
\end{equation}

The corresponding wavelength:
\begin{equation}
\lambda_{\rm cross} = \frac{2\pi}{k_{\rm cross}} \approx 3.7\times 10^{-34}~\text{m} \sim 23\,\ell_P.
\end{equation}

The corresponding frequency:
\begin{equation}
f_{\rm cross} = \frac{ck_{\rm cross}}{2\pi} \approx 8\times 10^{41}~\text{Hz} \sim \frac{1}{20\tau} \approx \frac{\gamma}{40}.
\end{equation}

\subsubsection{F2.3 Physical interpretation}

\begin{itemize}
\item \textbf{For $\lambda \gg \lambda_{\rm cross}$:} Overdamped diffusion. All astrophysical and laboratory scales (even LIGO gravitational waves at $\lambda \sim 10^5$~m) are in this regime. The parabolic $q$-field equation is an excellent approximation.
\item \textbf{For $\lambda \sim \lambda_{\rm cross}$:} Damped waves. This is the Planck scale---relevant only in the very early universe or for Planck-energy phenomena.
\item \textbf{For $\lambda \ll \lambda_{\rm cross}$:} Propagating waves with finite speed $c$. This regime may not be physically accessible for $q$-field perturbations (the cell scale $a \sim \ell_P$ imposes a UV cutoff), but it is mathematically present in the continuum hyperbolic extension.
\end{itemize}

\subsubsection{F2.4 Implication for gravitational waves}

LIGO observes gravitational waves at $f \sim 10$--$1000$~Hz, corresponding to $k \sim 10^{-6}$--$10^{-4}$~m$^{-1}$. This is $10^{40}$ times smaller than $k_{\rm cross}$. Gravitational waves are thus \textbf{deep in the overdamped regime}, meaning:
\begin{itemize}
\item The parabolic approximation is adequate for computing the static gravitational field.
\item However, gravitational waves \textbf{require the wave equation} to describe their propagation---the overdamped limit of the telegraph equation is still a wave equation (it has finite propagation speed).
\item The key point: $k \ll k_{\rm cross}$ means the $\partial_t^2$ term is negligible compared to the $\gamma\partial_t$ term for determining the dispersion, but it does NOT mean the propagation speed becomes infinite---the speed is always $c$, regardless of $k$.
\end{itemize}

\textbf{Important clarification:} The crossover $k_{\rm cross}$ separates regimes where the eigenvalues are real (overdamped) vs.\ complex (underdamped). It does NOT separate finite-speed from infinite-speed propagation. The Cattaneo system has finite speed $c$ for \textbf{all} $k$, because it is hyperbolic. The overdamped nature at small $k$ means perturbations decay monotonically without oscillation, but they still propagate at speed $\le c$.

% ============================================================
\subsection{F3. Planck-Scale Parameters: $a=\ell_P$, $\kappa = c^2/\ell_P^2$}
% ============================================================

\subsubsection{F3.1 Dimensional analysis}

The DGF framework introduces three dimensionful parameters:
\begin{itemize}
\item $a$: cell spacing [L]
\item $\tau$: cell update time [T]
\item $\kappa$: continuum diffusivity [$L^2$/T]
\end{itemize}

These are related by $\kappa \sim a^2/\tau$ (diffusivity from random walk scaling). The physical constants $c$, $G$, $\hbar$ must emerge from combinations of the DGF parameters.

\subsubsection{F3.2 The Planck system}

Define the Planck units:
\begin{align}
\ell_P &= \sqrt{\frac{\hbar G}{c^3}} \approx 1.616\times 10^{-35}~\text{m},\\
t_P &= \sqrt{\frac{\hbar G}{c^5}} \approx 5.391\times 10^{-44}~\text{s},\\
m_P &= \sqrt{\frac{\hbar c}{G}} \approx 2.176\times 10^{-8}~\text{kg}.
\end{align}

\subsubsection{F3.3 DGF parameter identification}

We identify:
\begin{align}
a &= \ell_P \quad \text{(cell spacing = Planck length)},\\
\tau &= t_P \quad \text{(cell time = Planck time)}.
\end{align}

Then the continuum diffusivity is:
\begin{equation}
\kappa = \frac{a^2}{\tau} = \frac{\ell_P^2}{t_P} = \ell_P c \approx 4.8\times 10^{-27}~\text{m}^2/\text{s}.
\label{eq:kappa_planck}
\end{equation}

\subsubsection{F3.4 Checking the Newtonian calibration}

The static calibration of the main paper requires $\kappa = 4\pi G/c^2$ (to match the Newtonian Poisson equation). Let us check whether (\ref{eq:kappa_planck}) is consistent:
\begin{equation}
\frac{4\pi G}{c^2} = \frac{4\pi \times 6.674\times 10^{-11}}{(2.998\times 10^8)^2} = 9.33\times 10^{-27}~\text{m}^2/\text{s}.
\end{equation}

The ratio:
\begin{equation}
\frac{4\pi G/c^2}{\ell_P c} = \frac{4\pi G}{c^2}\cdot\frac{c}{\sqrt{\hbar G/c}} = \frac{4\pi\sqrt{G}}{\sqrt{\hbar c}} = 4\pi\sqrt{\frac{G}{\hbar c}} = \frac{4\pi}{m_P c},
\end{equation}
which is not a pure number but has dimensions of 1/mass. This means $\kappa$ calibrated to Newtonian gravity and $\kappa_{\rm Planck} = \ell_P c$ differ by a dimensionful factor.

\textbf{Resolution:} Introduce a dimensionless constant $\alpha$ (the ``information-gravity coupling''):
\begin{equation}
\kappa = \alpha\,\ell_P c, \qquad \alpha = \frac{4\pi G/c^2}{\sqrt{\hbar G/c^3}\cdot c} = \frac{4\pi G}{c^3}\sqrt{\frac{c^3}{\hbar G}} = 4\pi\sqrt{\frac{G}{\hbar c}} \approx 1.94.
\end{equation}

The factor $\alpha \approx 1.94 \approx 2$ means $\kappa \approx 2\ell_P c$. The $\mathcal{O}(1)$ numerical factor is not predicted; its value must be determined by matching to experiment (which is what calibrating $\kappa$ to $G$ does).

\subsubsection{F3.5 The $\hbar$ dependence}

The appearance of $\hbar$ in $\ell_P$ means the DGF cell scale involves quantum mechanics. This is expected: the cell represents a qubit (quantum bit), and its physical realization must involve $\hbar$ (the quantum of action). However, \textbf{DGF in its current form does not incorporate quantum mechanics}---the dynamics $f:X\to X$ is classical (permutation on bit-strings).

\textbf{Tension:} The Planck-scale identification brings $\hbar$ into the theory, but the theory does not explain $\hbar$. This is an honest limitation. A fully quantum DGF (with cells as qubits and dynamics as unitary transformations on a Hilbert space of dimension $2^N$) would naturally produce $\hbar$ as the conversion factor between information and action, but this quantum extension does not yet exist.

\subsubsection{F3.6 Summary of the Planck-scale identification}

\begin{center}
\begin{tabular}{lll}
\toprule
DGF Parameter & Value & Relation \\
\midrule
$a$ (cell spacing) & $\ell_P$ & Input (identified with Planck length) \\
$\tau$ (cell time) & $t_P$ & Input ($=a/c$) \\
$\kappa$ (diffusivity) & $\alpha\ell_P c \approx 2\ell_P c$ & Calibrated ($\alpha$ from matching to $G$) \\
$\gamma = 1/\tau$ & $1/t_P \approx 1.85\times 10^{43}$~s$^{-1}$ & Derived from $\tau = t_P$ \\
$c$ & $c$ & Input (propagation speed bound from A3) \\
$G$ & --- & Calibrated: $G = \alpha\ell_P c^3/(4\pi)$ \\
$\hbar$ & --- & Enters via $\ell_P = \sqrt{\hbar G/c^3}$; not independently predicted \\
\bottomrule
\end{tabular}
\end{center}

\textbf{The bottom line:} If $a=\ell_P$ and $\tau=t_P$, then $\kappa$ and $G$ are related by $\kappa = \alpha\ell_P c$ and $G = \alpha\ell_P c^3/(4\pi)$. This single relation links two observed constants ($G$, $c$) to one fundamental length ($\ell_P$) and one dimensionless coupling ($\alpha$). The value $\alpha \approx 2$ is a numerical coincidence that the framework does not explain.

%%%%%%%%%%%%%%%%%%%%%%%%%%%%%%%%%%%%%%%%%%%%%%%%%%%%%%%%%%%%%%%%%%%%%
\section*{References for Supplemental Material}
%%%%%%%%%%%%%%%%%%%%%%%%%%%%%%%%%%%%%%%%%%%%%%%%%%%%%%%%%%%%%%%%%%%%%

\begin{thebibliography}{99}

\bibitem{Cattaneo1948} C.~Cattaneo, Atti Sem.\ Mat.\ Fis.\ Univ.\ Modena \textbf{3}, 83 (1948).
\bibitem{Vernotte1958} P.~Vernotte, C.\ R.\ Acad.\ Sci.\ \textbf{246}, 3154 (1958).
\bibitem{Joseph1989} D.D.\ Joseph and L.\ Preziosi, Rev.\ Mod.\ Phys.\ \textbf{61}, 41 (1989).
\bibitem{Jackson1971} H.E.\ Jackson, C.T.\ Walker, and T.F.\ McNelly, Phys.\ Rev.\ Lett.\ \textbf{25}, 26 (1970).
\bibitem{London1935} F.\ London and H.\ London, Proc.\ R.\ Soc.\ A \textbf{149}, 71 (1935).
\bibitem{Maxwell1865} J.C.\ Maxwell, Phil.\ Trans.\ R.\ Soc.\ \textbf{155}, 459 (1865).
\bibitem{Heaviside1876} O.\ Heaviside, Phil.\ Mag.\ \textbf{2}, 135 (1876).
\bibitem{Goldstein1951} S.~Goldstein, Quart.\ J.\ Mech.\ Appl.\ Math.\ \textbf{4}, 129 (1951).
\bibitem{Weiss1994} G.H.\ Weiss, \emph{Aspects and Applications of the Random Walk} (North-Holland, 1994).
\bibitem{Caldeira1983} A.O.\ Caldeira and A.J.\ Leggett, Ann.\ Phys.\ \textbf{149}, 374 (1983).
\bibitem{Will2014} C.M.\ Will, Living Rev.\ Relativ.\ \textbf{17}, 4 (2014).
\bibitem{LIGO2017} B.P.\ Abbott et al.\ (LIGO/Virgo), Astrophys.\ J.\ \textbf{848}, L13 (2017).
\bibitem{LIGO2019} B.P.\ Abbott et al.\ (LIGO/Virgo), Phys.\ Rev.\ Lett.\ \textbf{123}, 011102 (2019).
\bibitem{EHT2019} K.~Akiyama et al.\ (EHT), Astrophys.\ J.\ \textbf{875}, L1 (2019).
\bibitem{Shannon1948} C.E.\ Shannon, Bell Syst.\ Tech.\ J.\ \textbf{27}, 379 (1948).
\bibitem{Vazquez2007} J.L.\ Vazquez, \emph{The Porous Medium Equation} (Oxford, 2007).
\bibitem{Bak1987} P.~Bak, C.~Tang, and K.~Wiesenfeld, Phys.\ Rev.\ Lett.\ \textbf{59}, 381 (1987).
\bibitem{Zurek2009} W.H.\ Zurek, Nat.\ Phys.\ \textbf{5}, 181 (2009).

\end{thebibliography}

\end{document}
